\chapter{相对论动力学与电磁场}
\section{相位与波四矢量}\label{th:m3wave}
单色平面波的复表示为$\mathcal E=\mathcal E_0\exp[i(\vct k\cdot\vct r-\omega t)]$。同一波面的事件在不同坐标系中仍具有相同相位，所以$\omega t-\vct k\cdot\vct r$是洛伦兹标量。将逆坐标变换代入相位，按$t',x',y',z'$收集系数，得到
\begin{align}\label{eq:m3wave}
 \omega'&=\gamma(\omega-Vk_x),& k_x'&=\gamma(k_x-V\omega/c^2),\\
 k_y'&=k_y,&k_z'&=k_z.\notag
\end{align}
因此$k^\mu=(\omega/c,\vct k)$是波四矢量。真空中$\omega=c|\vct k|$，故$k^2=0$，任何惯性系中仍有真空色散关系。若在介质中使用波矢变换，介质的运动也要一起变换，不能同时假定两系介质都静止且折射率相同。

由$k_x=(\omega/c)\cos\theta$可得$\omega'=\gamma\omega(1-\beta\cos\theta)$。再以$k_x'c/\omega'$求方向，便重现光行差公式。频率和方向应作为同一四矢量的分量一起求出。

\begin{example}\label{ex:m3movingmirror}
一束频率$\nu$的光沿$x$正向照射以速度$V>0$远离的理想平面镜，镜面垂直$x$。忽略镜子反冲。求反射光在实验室中的频率。
\Diagram{fig:m3mirrorq}{运动镜面与入射、反射光}{\draw[curve](4,-1)--(4,1.4);\draw[axis](4.2,1)--(5.5,1)node[above]{$V$};\draw[very thick,->](0,.6)--(3.7,.6)node[midway,above]{$\nu$};\draw[very thick,->](3.7,-.3)--(.2,-.3)node[midway,below]{$\nu_r$};}
\end{example}
\begin{solution}
先把图\ref{fig:m3mirrorq}的入射光变到镜系。入射$k_x=\omega/c$，故$\omega'_i=\gamma\omega(1-\beta)$。镜系中镜面静止，反射保持频率并使$k_x'$反号：$\omega'_r=\omega'_i$，$k'_{xr}=-\omega'_i/c$。
\Diagram{fig:m3mirrora}{镜系反射保持频率，再将反向波变回实验室}{\draw[curve](3,-1)--(3,1.3);\draw[->](0,.5)--(2.8,.5)node[midway,above]{$\omega'_i$};\draw[->](2.8,-.4)--(0,-.4)node[midway,below]{$\omega'_r=\omega'_i$};\node[right]at(3.2,.2){镜系};}
逆变换$\omega_r=\gamma(\omega'_r+Vk'_{xr})$给$\omega_r=\gamma(1-\beta)\omega'_i$。因此图\ref{fig:m3mirrora}对应的两次频率变换合成
\[
 \frac{\nu_r}\nu=\gamma^2(1-\beta)^2=\frac{1-\beta}{1+\beta}.
\]
若镜子靠近，只需令$\beta<0$，频率升高。低速时$\nu_r/\nu\simeq1-2\beta$，因频移发生在入射与返回两段，不是一次多普勒因子。
\end{solution}

\section{总四动量、质心系与阈值}\label{th:m3momentum}
封闭系统的总四动量$P^\mu=\sum_i p_i^\mu$守恒。若$P^2>0$，存在总空间动量为零的质心系，系统不变质量为
\[
 M^2c^4=E_{\rm tot}^2-c^2\vct P^2.
\]
其速度为$\vct V_{\rm cm}=c^2\vct P/E_{\rm tot}$。在该系总能量等于$Mc^2$，其中包含各粒子的相对运动能、内部能及相互作用能；一般$M\ne\sum_i m_i$。讨论散射时取入射和出射粒子相隔很远的状态，可把相互作用能归入粒子静质量及可分辨辐射，直接使用粒子四动量求和。

\subsection{功—能关系与一维加速}
单粒子满足$E^2=p^2c^2+m^2c^4$，故$\dd E=c^2\vct p\cdot\dd\vct p/E=\vct v\cdot\dd\vct p$。由此$\dd E=\vct f\cdot\dd\vct r$，无须先假定牛顿动能。对直线运动，$\dd p/\dd v=m\gamma_v^3$；对瞬时垂直方向的速度微扰，$\dd\vct p_\perp=m\gamma_v\dd\vct v_\perp$。相同三力在两方向产生的实验室加速度不同。

\begin{example}\label{ex:m3photonpair}
两个光子能量为$E_1,E_2$，传播方向夹角$\theta$。求它们组成系统的不变质量；讨论平行、反向两种极限及能否产生静质量$m$的正反粒子对。
\Diagram{fig:m3pairq}{两光子四动量的空间夹角}{\draw[axis](0,0)--(3.6,0)node[right]{$\vct p_1$};\draw[axis](0,0)--(1.8,2.5)node[right]{$\vct p_2$};\draw(.8,0)arc(0:54.25:.8);\node at(1,.5){$\theta$};}
\end{example}
\begin{solution}
图\ref{fig:m3pairq}中$|\vct p_i|=E_i/c$，所以
\begin{align*}
 M^2c^4&=(E_1+E_2)^2-c^2(p_1^2+p_2^2+2p_1p_2\cos\theta)\\
 &=2E_1E_2(1-\cos\theta).
\end{align*}
平行时$M=0$，不存在质心静止系；反向时$Mc^2=2\sqrt{E_1E_2}$。若反向且$E_1=E_2$，实验室就是质心系。
\Diagram{fig:m3paira}{质心系中的粒子对阈值：末态相对动量趋于零}{\draw[axis](-2.7,0)--(2.7,0);\draw[->,thick](-2.5,.6)--(-.3,.6)node[midway,above]{$E'_1$};\draw[->,thick](2.5,.6)--(.3,.6)node[midway,above]{$E'_2$};\fill(-.5,-.5)circle(2pt)node[below]{$m$};\fill(.5,-.5)circle(2pt)node[below]{$m$};}
图\ref{fig:m3paira}所示的最小末态质心能量为$2mc^2$，因而阈值条件为$E_1E_2(1-\cos\theta)\geq2m^2c^4$。这只是能量动量允许的条件，还须满足电荷等其他守恒律和实际相互作用机制。单个自由光子不可能在真空中产生一对有质量粒子，因为其总不变质量为零。
\end{solution}

\begin{example}\label{ex:m3inelastic}
静质量$m$的粒子以动能$K$撞击同样质量的静止粒子，两者完全粘合且不向外辐射。求粘合体的静质量、速度及其内部增加的能量。
\Diagram{fig:m3inelasticq}{实验室中的完全非弹性碰撞}{\fill(0,0)circle(3pt)node[below]{$m$};\draw[axis](.2,.4)--(1.8,.4)node[above]{$K$};\fill(3,0)circle(3pt)node[below]{$m$静止};\draw[->](3.5,0)--(4.5,0);\draw[curve](5,0)circle(.35);\node[below]at(5,-.35){$M$};}
\end{example}
\begin{solution}
图\ref{fig:m3inelasticq}中，入射粒子的总能量为$E_1=mc^2+K$，靶粒子能量为$mc^2$。入射动量由质壳关系给出：
\[
 p^2c^2=E_1^2-m^2c^4=K^2+2mc^2K.
\]
粘合体携带全部四动量，所以$E_T=2mc^2+K$、$P_T=p$，其静质量满足
\begin{align*}
 M^2c^4&=(2mc^2+K)^2-p^2c^2\\
 &=4m^2c^4+2mc^2K.
\end{align*}
故$M=2m\sqrt{1+K/(2mc^2)}$，速度为$V=Pc^2/E_T=c\sqrt{K^2+2mc^2K}/(2mc^2+K)$。
\Diagram{fig:m3inelastica}{质心系内相向运动的动能变为粘合体内部能}{\draw[->](-2,.4)--(-.5,.4);\draw[->](2,.4)--(.5,.4);\fill(-1.5,0)circle(2.5pt);\fill(1.5,0)circle(2.5pt);\draw[curve](0,-1)circle(.35);\node[right]at(.5,-1){$\Delta U=(M-2m)c^2$};}
图\ref{fig:m3inelastica}是在总动量为零的质心系，粘合前的相向动能全部转为内部能，因此
\[
 \Delta U=(M-2m)c^2,
 \qquad K_{\rm bulk}=E_T-Mc^2=K-\Delta U.
\]
两项之和为初始实验室动能$K$。其中$\Delta U$是粘合体自身系的内能增加，$K_{\rm bulk}$是实验室中的整体动能。令$k=K/(2mc^2)$，展开$\sqrt{1+k}=1+k/2-k^2/8+\cdots$，得到
\[
 \Delta U=\frac K2-\frac{K^2}{16mc^2}+\cdots,\qquad
 K_{\rm bulk}=\frac K2+\frac{K^2}{16mc^2}+\cdots.
\]
牛顿极限中两者各占一半。相对论情形不能把所有入射动能都算成内部能，因为总动量守恒要求粘合体继续平动。
\end{solution}

\section{四力、加速度与瞬时静止系}\label{th:m3force}
设实验室粒子速度为$\vct u$，使用$\gamma_u$区别于参考系相对运动的$\gamma$。三力$\vct f=\dd\vct p/\dd t$对应
\[
 F^\mu=\gamma_u(\vct f\cdot\vct u/c,f_x,f_y,f_z).
\]
沿$x$变换后，粒子的新洛伦兹因子$\gamma_{u'}=\gamma\gamma_u(1-Vu_x/c^2)$。以$F'^i/\gamma_{u'}$求三力，得
\begin{align}\label{eq:m3forcetrans}
 f'_x&=\frac{f_x-V(\vct f\cdot\vct u)/c^2}{1-Vu_x/c^2},\\
 f'_{y,z}&=\frac{f_{y,z}}{\gamma(1-Vu_x/c^2)}.\notag
\end{align}
在瞬时静止系中$u'=0$，四力时间分量为零，$F^2=-f_0^2$。这里$f_0$专指瞬时静止系三力的大小，固定惯性系中的三力须另作变换。

四加速度需同时求导$\gamma_u$与$\vct u$：
\[
 \dot\gamma_u=-\frac12(1-u^2/c^2)^{-3/2}
 \left(-\frac{2\vct u\cdot\vct a}{c^2}\right)
 =\gamma_u^3\frac{\vct u\cdot\vct a}{c^2}.
\]
再用$\dd/\dd\tau=\gamma_u\dd/\dd t$对$u^\mu=\gamma_u(c,\vct u)$逐项求导，得
\[
 a^0=\gamma_u^4\frac{\vct u\cdot\vct a}{c},\quad
 \vct a_{(4)}=\gamma_u^2\vct a+\gamma_u^4\frac{\vct u\cdot\vct a}{c^2}\vct u.
\]
沿瞬时速度方向分解，四加速度的纵、横空间分量分别为$\gamma_u^4a_\parallel$、$\gamma_u^2a_\perp$，时间分量为$\gamma_u^4ua_\parallel/c$。因此
\[
 \alpha^2=-a_\mu a^\mu
 =\gamma_u^8(1-u^2/c^2)a_\parallel^2+\gamma_u^4a_\perp^2
 =\gamma_u^6a_\parallel^2+\gamma_u^4a_\perp^2.
\]
恒静质量时$F^\mu=ma^\mu$，这也给出四力模与固有加速度的关系。

\begin{example}\label{ex:m3accel}
粒子在实验室瞬时以速度$u$沿$x$运动，其瞬时静止系加速度分量为$(\alpha_x,\alpha_y,0)$。从四加速度变换推导实验室两分量。
\Diagram{fig:m3accelq}{瞬时静止系中的加速度分解}{\draw[axis](0,0)--(3.8,0)node[right]{$x_0$};\draw[axis](0,0)--(0,2.8)node[above]{$y_0$};\draw[->,thick](0,0)--(2.4,1.8);\draw[guide](2.4,0)--(2.4,1.8)--(0,1.8);\node[below]at(1.2,0){$\alpha_x$};\node[left]at(0,.9){$\alpha_y$};}
\end{example}
\begin{solution}
把四加速度记为$A^\mu=\dd U^\mu/\dd\tau$，实验室三加速度记为$\vct a$，以免同一字母混淆。图\ref{fig:m3accelq}的瞬时静止系中，$A_*^\mu=(0,\alpha_x,\alpha_y,0)$。沿$x$逆变换得到
\[
 A^0=\gamma_u\beta_u\alpha_x,\quad
 A^x=\gamma_u\alpha_x,\quad A^y=\alpha_y.
\]
这里变换的是同一事件处的四矢量，不把该瞬时静止系延伸为整个加速过程的静止系。
\Diagram{fig:m3accela}{实验室加速度的纵、横分量分别多出不同的时间因子}{\draw[axis](0,0)--(3.8,0)node[right]{$x,\ u$};\draw[axis](0,0)--(0,2.8)node[above]{$y$};\draw[->,thick](0,0)--(1.3,1.8);\draw[guide](1.3,0)--(1.3,1.8)--(0,1.8);\node[below]at(1.3,0){$\alpha_x/\gamma_u^3$};\node[left]at(0,1){$\alpha_y/\gamma_u^2$};}
为将图\ref{fig:m3accela}中的$A^i$还原为三加速度，先对$\gamma_u$求导：
\[
 \frac{\dd\gamma_u}{\dd t}=\gamma_u^3\frac{\vct u\cdot\vct a}{c^2}
 =\gamma_u^3\frac{ua_x}{c^2}.
\]
因$\dd/\dd\tau=\gamma_u\dd/\dd t$，该事件处$u_y=0$，所以
\begin{align*}
 A^x&=\gamma_u\frac{\dd(\gamma_u u_x)}{\dd t}
 =\gamma_u^2a_x+\gamma_u^4\frac{u^2a_x}{c^2}
 =\gamma_u^4a_x,\\
 A^y&=\gamma_u\frac{\dd(\gamma_u u_y)}{\dd t}=\gamma_u^2a_y.
\end{align*}
与变换所得分量比较，得到$a_x=\alpha_x/\gamma_u^3$、$a_y=\alpha_y/\gamma_u^2$。时间分量$A^0=\gamma_u^4ua_x/c$也与上式一致。纵向多出的一个因子来自速度变化引起的$\gamma_u$变化，不能把横向时间膨胀的推算直接套到纵向。
\end{solution}

\begin{example}\label{ex:m3longfield}
均匀电场沿$x$，带电粒子瞬时也沿$x$运动。用瞬时静止系三力和加速度关系说明为何平行电场不变；再求实验室加速度。
\Diagram{fig:m3longq}{平行电场中的带电粒子}{\foreach\y in{0,.6,1.2}{\draw[axis](0,\y)--(4,\y);}\node[right]at(4,.6){$E_x$};\fill(1.5,.6)circle(3pt)node[above]{$q,m$};\draw[axis](1.7,1.6)--(3,1.6)node[above]{$u$};}
\end{example}
\begin{solution}
图\ref{fig:m3longq}中实验室$f_x=qE_x$。把式\eqref{eq:m3forcetrans}取$V=u$且$\vct f\cdot\vct u=f_xu$，得$f_{0x}=f_x$。粒子在瞬时静止系受到电力$qE_{0x}$，故$qE_{0x}=qE_x$，即$E_{0x}=E_x$。
\Diagram{fig:m3longa}{相同平行电力对应不同的坐标加速度}{\draw[axis](0,0)--(4.3,0)node[right]{$u/c$};\draw[axis](0,0)--(0,2.8)node[above]{$a_x/(qE_x/m)$};\draw[curve,domain=0:1,samples=70]plot({3.5*\x},{2.3*(1-\x*\x)^1.5});\node[left]at(0,2.3){$1$};\node[below]at(3.5,0){$1$};}
瞬时静止系$\alpha_x=qE_x/m$，由例\ref{ex:m3accel}，图\ref{fig:m3longa}的实验室加速度$a_x=qE_x/(m\gamma_u^3)$。随着速度增加，电力保持不变，加速度减小；能量输入率仍为$qE_xu$，并不因加速度变小而趋于零。
\end{solution}

\section{四电流与电荷守恒}\label{th:m3current}
对于共同运动的带电粒子群，静止系密度为$\rho_0$，实验室电荷密度$\rho=\gamma_u\rho_0$，电流$\vct j=\rho\vct u$，因此
\[
 J^\mu=(c\rho,\vct j)=\rho_0u^\mu,\quad J^2=c^2\rho_0^2.
\]
多种不同速度的带电成分应分别构造四电流再求和；总四电流未必类时，也未必存在把所有电荷同时静止的惯性系。连续性方程
\[
 \frac{\partial\rho}{\partial t}+\nabla\cdot\vct j=0
\]
可以写为$\partial_\mu J^\mu=0$，其积分形式是某固定体积内电荷减少等于流出边界的电荷。沿$x$的变换给$\rho'=\gamma(\rho-Vj_x/c^2)$、$j'_x=\gamma(j_x-V\rho)$，横向电流分量不变。

\begin{example}\label{ex:m3wire}
无限直导线沿$x$，实验室中正离子静止，线电荷密度$+\lambda$；电子速度为$-u$，电子实验室线密度为$-\lambda$。求沿$+V$运动的观察系中的总线电荷密度和电流。规定常规电流沿正向。
\Diagram{fig:m3wireq}{实验室中电中性的载流导线}{\draw[curve](-.3,0)--(5.8,0);\foreach\x in{0,1,2,3,4,5}{\node at(\x,.4){$+$};\node at(\x,-.4){$-$};}\draw[axis](3,-.9)--(1.5,-.9)node[midway,below]{$u$};\draw[axis](1,1)--(2.5,1)node[midway,above]{$I=\lambda u$};}
\end{example}
\begin{solution}
图\ref{fig:m3wireq}的总$\lambda_T=0$，$I=\lambda u$。对单位横截面的积分四电流同样变换，得到
\[
 \lambda_T'=-\gamma\frac{VI}{c^2},\qquad I'=\gamma I.
\]
也可逐成分求解：离子$\lambda'_+=\gamma\lambda$；电子$\lambda'_- =\gamma[-\lambda-V(\lambda u)/c^2]$。两者相加即得上式。
\Diagram{fig:m3wirea}{运动观察系中负电荷密度的绝对值大于正电荷密度}{\draw[curve](0,0)--(5.8,0);\foreach\x in{.3,1.3,2.3,3.3,4.3,5.3}{\node at(\x,.4){$+$};}\foreach\x in{0,.7,1.4,2.1,2.8,3.5,4.2,4.9,5.6}{\node at(\x,-.4){$-$};}\node[below]at(3,-1){$\lambda'_T<0$，示意密度};}
图\ref{fig:m3wirea}说明电中性是与参考系有关的密度条件，整段材料上包含多少总电荷则还涉及边界和同时截面。若直接把离子和电子都缩短相同倍数，便错误地忽略了电子在实验室已经运动以及密度变换中的电流项。
\end{solution}

\section{四势、规范与场方程}\label{th:m3potential}
由$\nabla\cdot\vct B=0$，局部可写$\vct B=\nabla\times\vct A$。法拉第定律给$\nabla\times(\vct E+\partial_t\vct A)=0$，所以$\vct E=-\nabla\phi-\partial_t\vct A$。势不唯一：$\vct A\to\vct A+\nabla\zeta$、$\phi\to\phi-\partial_t\zeta$保持场不变。

把势代入高斯定律，得到$-\nabla^2\phi-\partial_t\nabla\cdot\vct A=\rho/\varepsilon_0$。取洛伦兹规范
\[
 \nabla\cdot\vct A+\frac1{c^2}\partial_t\phi=0,
\]
就有$(c^{-2}\partial_t^2-\nabla^2)\phi=\rho/\varepsilon_0$。再把势代入安培—麦克斯韦定律，利用$\nabla\times\nabla\times\vct A=\nabla(\nabla\cdot\vct A)-\nabla^2\vct A$，梯度项由规范条件抵消，得到
\[
 \left(\frac1{c^2}\partial_t^2-\nabla^2\right)\vct A=\mu_0\vct j.
\]
令$A^\mu=(\phi/c,\vct A)$、$\Box=c^{-2}\partial_t^2-\nabla^2$，两式合并为$\Box A^\mu=\mu_0J^\mu$。在保持此协变规范的选择下，四势按四矢量变换。若另行改用其他规范，还须补相应规范变换，不能把某个规范下的分量形式当作场本身。

\subsection{导数变换}
由链式法则，$\partial_{ct'}=\gamma(\partial_{ct}+\beta\partial_x)$、$\partial_{x'}=\gamma(\partial_x+\beta\partial_{ct})$，$\partial_{y',z'}=\partial_{y,z}$。特别注意导数的正负号与坐标正变换不同，因为这里是对逆坐标式求导。将两导数组合平方，交叉项抵消，就可验证$\Box'=\Box$。

\section{电磁场的完整变换}\label{th:m3fields}
对四势作沿$x$的变换，$\phi'=\gamma(\phi-VA_x)$、$A_x'=\gamma(A_x-V\phi/c^2)$、$A_y'=A_y$、$A_z'=A_z$。例如
\begin{align*}
 E_y'&=-\partial_y\phi'-\partial_{t'}A_y'\\
 &=-\gamma\partial_y\phi+\gamma V\partial_y A_x
 -\gamma\partial_tA_y-\gamma V\partial_xA_y\\
 &=\gamma(E_y-VB_z).
\end{align*}
这里$B_z=\partial_xA_y-\partial_yA_x$，负号由叉乘方向决定。同样将$E_x'$完整展开，时间与空间交叉导数抵消，剩下$\gamma^2(1-\beta^2)E_x=E_x$。对$\vct B'=\nabla'\times\vct A'$逐分量计算，得到
\begin{equation}\label{eq:m3field}
 \begin{aligned}
 E_x'&=E_x,&B_x'&=B_x,\\
 E_y'&=\gamma(E_y-VB_z),&B_y'&=\gamma(B_y+VE_z/c^2),\\
 E_z'&=\gamma(E_z+VB_y),&B_z'&=\gamma(B_z-VE_y/c^2).
 \end{aligned}
\end{equation}
相对任意$\vct V$的矢量形式是
\[
 \vct E'_\parallel=\vct E_\parallel,\quad
 \vct E'_\perp=\gamma(\vct E+\vct V\times\vct B)_\perp,
\]
\[
 \vct B'_\parallel=\vct B_\parallel,\quad
 \vct B'_\perp=\gamma(\vct B-\vct V\times\vct E/c^2)_\perp.
\]
从分量式可直接验证$E^2-c^2B^2$和$\vct E\cdot\vct B$不变。两个不变量可以排除某些变换目标，但具体参考系速度仍须代入变换式求出。

\begin{example}\label{ex:m3crossed}
实验室中$\vct E=E\hat{\vct y}$、$\vct B=B\hat{\vct z}$，$E,B>0$。求能使电场或磁场消失的沿$x$运动惯性系，并讨论$E=cB$。
\Diagram{fig:m3crossq}{正交电磁场与可选的变换方向}{\draw[axis](0,0)--(3.6,0)node[right]{$x$};\draw[axis](0,0)--(0,2.5)node[above]{$E\hat{\vct y}$};\draw(1.5,1)circle(.18);\fill(1.5,1)circle(1.2pt);\node[right]at(1.8,1){$B\hat{\vct z}$};}
\end{example}
\begin{solution}
图\ref{fig:m3crossq}中$E_y'=\gamma(E-VB)$、$B_z'=\gamma(B-VE/c^2)$。若$E<cB$，取$V=E/B<c$，可令$E'=0$。此时$B'=\sqrt{B^2-E^2/c^2}$，仍沿$z$正向。
\Diagram{fig:m3crossa}{由不变量符号区分磁占优、电占优和光状场}{\draw[axis](0,0)--(5,0)node[right]{$E/(cB)$};\draw[curve](2.5,-.6)--(2.5,.8);\node[below]at(2.5,-.6){$1$};\node[above]at(1,.2){可令$E'=0$};\node[above]at(4,.2){可令$B'=0$};}
若$E>cB$，取$V=c^2B/E<c$，得到$B'=0$和$E'=\sqrt{E^2-c^2B^2}$。图\ref{fig:m3crossa}的分界$E=cB$要求$V=c$才能消场，不是合法惯性系。此时两个不变量均为零，任何有限速度系中仍有非零正交场且$E'=cB'$，真空平面波正属于这类情况。
\end{solution}

\begin{example}\label{ex:m3wavefield}
真空平面波沿$x$传播，电场沿$y$，磁场沿$z$。沿波方向以速度$V$运动的观察者测得何种频率、振幅与光强？
\Diagram{fig:m3waveq}{沿波传播方向追赶平面波}{\draw[axis](-.2,0)--(5.5,0)node[right]{$\vct k,V$};\draw[curve,domain=0:5,samples=100]plot(\x,{.7*sin(180*\x)});\node[above]at(2,1){$E_y$};\draw[axis](3,-1)--(4.5,-1);}
\end{example}
\begin{solution}
图\ref{fig:m3waveq}中$B_z=E_y/c$。式\eqref{eq:m3field}给$E_y'=\gamma(1-\beta)E_y$、$B_z'=E_y'/c$。频率同样有$\omega'=\gamma(1-\beta)\omega$。
\Diagram{fig:m3wavea}{振幅和频率同比例红移，波长相应伸长}{\draw[axis](0,0)--(5.6,0)node[right]{$x'$};\draw[curve,domain=0:5,samples=90]plot(\x,{.45*sin(100*\x)});\draw[<->](.9,-.9)--(4.5,-.9)node[midway,below]{$\lambda'=\lambda/[\gamma(1-\beta)]$};}
时间平均光强$I=\varepsilon_0cE_0^2/2$，所以图\ref{fig:m3wavea}中的$I'/I=\gamma^2(1-\beta)^2=(1-\beta)/(1+\beta)$。频率比和振幅比相同，而光强比是其平方。此结果对整个无限平面波的局部能流成立，与点源角功率的立体角雅可比属于不同问题。
\end{solution}
\section{守恒方程中的未知量与参考系}
\subsection{质心系的速度和能量}
总四动量为$(E/c,\vct P)$，若$E^2-c^2P^2>0$，令参考系速度沿$\vct P$，大小$V=c^2P/E$。变换后的平行动量为$P'=\gamma(P-VE/c^2)=0$。相应$\gamma=E/\sqrt{E^2-c^2P^2}$，新系能量
\[
 E'=\gamma(E-VP)=\sqrt{E^2-c^2P^2}=Mc^2.
\]
总四动量为类时时，质心系速度由它唯一确定。同向光子束的总四动量为类光，形式上需要$V=c$，因而没有质心惯性系。固定靶实验的实验室一般有非零总动量，阈值时末态仍须携带整体动能。

\subsection{两体末态的通用解}
在质心系中，两末态粒子的动量必为$\vct p_*,-\vct p_*$。设总能量为$W$，质量$m_1,m_2$，则$E_{1*}+E_{2*}=W$，并有$E_{1*}^2-E_{2*}^2=(m_1^2-m_2^2)c^4$。将左侧因式分解，得到
\[
 E_{1*}=\frac{W^2+(m_1^2-m_2^2)c^4}{2W},\quad
 E_{2*}=\frac{W^2+(m_2^2-m_1^2)c^4}{2W}.
\]
再由$E_{1*}^2-p_*^2c^2=m_1^2c^4$，得到
\[
 p_*c=\frac{\sqrt{[W^2-(m_1+m_2)^2c^4]
 [W^2-(m_1-m_2)^2c^4]}}{2W}.
\]
末态粒子各有正能量$E_{i*}=\sqrt{m_i^2c^4+p_*^2c^2}\geq m_ic^2$，所以物理分支必须满足$W\geq(m_1+m_2)c^2$，这才是两体阈值。仅要求根号为实还可能保留$W\leq|m_1-m_2|c^2$的代数分支，它不满足上述两个正能量条件。在阈值极限$p_*=0$，末态没有相对运动，但整体在实验室仍可运动。内部激发应计入$m_i$，不能在公式外再次添加同一激发能。

\Diagram{fig:m3twobody}{两体质心系中动量等大反向，能量未必相等}{\draw[axis](-3,0)--(3.2,0);\fill(-1,0)circle(2pt)node[below]{$m_1$};\fill(1,0)circle(2pt)node[below]{$m_2$};\draw[->,thick](-1,.5)--(-2.7,.5)node[above]{$p_*$};\draw[->,thick](1,.5)--(2.7,.5)node[above]{$p_*$};}
图\ref{fig:m3twobody}给出反向动量，不要求两个粒子以相同速率运动。若$m_1\ne m_2$，能量也不同。得到$E_{i*},p_*$后，再按题目给定出射角变回实验室，通常比直接在实验室同时联立角度、速度、能量更清楚。

\subsection{相对速度与入射能量不是同一输入}
在两个相同粒子的对撞中，某系观察到每束速率为$v$，彼此相对速率是$2v/(1+v^2/c^2)$。质心能量却是$2\gamma_vmc^2$，不能把相对速率直接代入“每束能量”公式。相对速率属于一粒子的静止系，束能量属于实验室，所选参考系必须在整个守恒式中一致。

同样，“完全非弹性”只规定末态相对运动消失，不规定静质量守恒。动能变为热或激发后，总四动量仍守恒，而复合物静质量增加。若相互作用放出辐射，须把辐射四动量留在末态，复合体的质量和速度会相应改变；不能在不封闭的子系统中同时要求能量、动量各自不变。

\section{四势推导中的全部横向分量}
\subsection{磁场变换的一项完整展开}
以$B_z'=\partial_{x'}A_y'-\partial_yA_x'$为例，代入导数和势的变换：
\begin{align*}
 B_z'&=\gamma(\partial_x+V\partial_t/c^2)A_y
 -\gamma\partial_y(A_x-V\phi/c^2)\\
 &=\gamma(\partial_xA_y-\partial_yA_x)
 +\frac{\gamma V}{c^2}(\partial_tA_y+\partial_y\phi)\\
 &=\gamma(B_z-VE_y/c^2).
\end{align*}
每个括号都是原系已定义的场，不需要假设场静态。对$B_y'=\partial_zA_x'-\partial_{x'}A_z'$，同样给$\gamma(B_y+VE_z/c^2)$；符号变化来自旋度的分量次序。$B_x'=\partial_yA_z'-\partial_zA_y'=B_x$，因为这两势分量和这两个导数都未变。

纵向电场中抵消较多，完整写成
\begin{align*}
 E_x'&=-\gamma^2(\partial_x+V\partial_t/c^2)(\phi-VA_x)\\
 &\quad-\gamma^2(\partial_t+V\partial_x)(A_x-V\phi/c^2).
\end{align*}
展开后，$V\partial_xA_x$和$V\partial_t\phi/c^2$的交叉项分别抵消，余下$-\gamma^2(1-\beta^2)(\partial_x\phi+\partial_tA_x)=E_x$。若漏掉$\partial_{t'}$中的空间导数项，就会错误地产生纵向场也乘$\gamma$的结果。

\subsection{场不变量的代数核对}
沿$x$推动时，纵向$E_x^2-c^2B_x^2$不变。将横向组合为$E_y'^2-c^2B_z'^2$，代入式\eqref{eq:m3field}：
\begin{align*}
 E_y'^2-c^2B_z'^2
 &=\gamma^2[(E_y-VB_z)^2-c^2(B_z-VE_y/c^2)^2]\\
 &=\gamma^2(1-\beta^2)(E_y^2-c^2B_z^2)\\
 &=E_y^2-c^2B_z^2.
\end{align*}
另一对$(E_z,B_y)$相同，三项相加就得$E'^2-c^2B'^2=E^2-c^2B^2$。点积中的横向交叉项也成对抵消，得到$\vct E'\cdot\vct B'=\vct E\cdot\vct B$。

非零真空平面波的两个场不变量均为零，其能量密度和坡印廷矢量仍非零。有限洛伦兹变换可逆，所以全零场只能由全零场变来；非零场至多在满足条件时消去其中一种场。

\section{麦克斯韦方程与电荷连续性}
安培—麦克斯韦方程为$\nabla\times\vct B=\mu_0\vct j+c^{-2}\partial_t\vct E$。取散度，左侧旋度的散度恒为零，右侧用$\nabla\cdot\vct E=\rho/\varepsilon_0$，得$0=\mu_0\nabla\cdot\vct j+\mu_0\partial_t\rho$。因此连续性方程是场方程相容的要求，位移电流项对变化电荷必不可少。

采用$\mathcal F^{0i}=-E_i/c$、$\mathcal F^{i0}=E_i/c$、$\mathcal F^{ij}=-\epsilon_{ijk}B_k$，其中$i,j,k=1,2,3$，完整矩阵及洛伦兹力的分量核对见例\ref{ex:m3fourtensor}。含源方程为$\partial_\nu\mathcal F^{\nu\mu}=\mu_0J^\mu$。取$\mu=0$，空间散度给$\nabla\cdot\vct E/c=\mu_0c\rho$，就是高斯定律；取空间分量，得到$\nabla\times\vct B-c^{-2}\partial_t\vct E=\mu_0\vct j$。再缩并$\partial_\mu$，由于$\partial_\mu\partial_\nu$对称而场张量反对称，左侧恒零，故$\partial_\mu J^\mu=0$。指标次序须与场张量定义一致。

\subsection{规范自由度与实际场}
洛伦兹规范下$\Box A^\mu=\mu_0J^\mu$，规范函数若满足$\Box\zeta=0$，改变势后仍保持同一规范。因此即使固定洛伦兹规范，势也未完全唯一。可观测场由势的导数组合决定，某个分量为零不能单独作为“没有电场”或“没有磁场”的证据。

对局域源，场方程的物理解还需选择边界条件，例如不从无限远无缘无故入射的辐射，以及对过去源历史的延迟响应。仅验证某个势满足波动方程，并不能说明它已满足实际初始、边界条件。匀速电荷场、加速电荷辐射场、外加平面波虽然都遵循同一方程，其边界数据不同。

\subsection{场的能量与参考系}
真空电磁能量密度为$u_{\rm em}=\varepsilon_0E^2/2+B^2/(2\mu_0)$，能流$\vct S=\vct E\times\vct B/\mu_0$。纯静电场的能流可以为零，变到运动系出现磁场后能流一般非零，这与能量和动量共同变换一致。不能只变换能量密度而保留原系“动量为零”的判断。

\Diagram{fig:m3energyflow}{运动电荷场的电磁能流与场的角分布相联系}{\draw[axis](-2.3,0)--(2.7,0)node[right]{$V$};\fill(0,0)circle(2pt);\draw[axis](0,1.5)--(0,2.5)node[above]{$E_y$};\draw(.5,1.6)circle(.14);\fill(.5,1.6)circle(1pt);\draw[axis](.6,2.3)--(2,2.3)node[right]{$\vct S$};}
图\ref{fig:m3energyflow}在电荷横向上方给$E_y$、$B_z$，所以能流沿正$x$。局部能量守恒为$\partial_tu_{\rm em}+\nabla\cdot\vct S=-\vct j\cdot\vct E$。其中$\vct j\cdot\vct E$是场对物质做功的功率密度，右侧的负号表示这部分能量从场中流失；磁力不直接做功。不同参考系可有不同的电力、磁力分量，总能量动量守恒仍成立。

\subsection{电磁场不变量怎样限制可选参考系}
若某系中电场为零，则$\vct E\cdot\vct B=0$且$c^2B^2-E^2>0$；若磁场为零，则第一不变量仍为零，而第二不变量小于零。这些是必要条件，也给出正交场中的构造方法。设$\vct E\perp\vct B$，沿$\vct E\times\vct B$方向取速度，场变换中平行速度的分量均为零，横向分量为
\[
 \vct E'=\gamma(\vct E+\vct V\times\vct B),\quad
 \vct B'=\gamma(\vct B-\vct V\times\vct E/c^2).
\]
磁占优时取$\vct V=\vct E\times\vct B/B^2$。矢量恒等式$(\vct E\times\vct B)\times\vct B=-B^2\vct E$使$\vct E'=0$，而$V=E/B<c$。电占优时取$\vct V=c^2\vct E\times\vct B/E^2$，得到$\vct B'=0$，且$V=c^2B/E<c$。

\Diagram{fig:m3fieldclasses}{正交场沿电磁能流方向选择变换；等强边界不能由有限惯性变换静止化}{\draw[axis](0,0)--(4,0)node[right]{$\vct E\times\vct B$};\draw[->,thick](0,0)--(0,2.5)node[above]{$\vct E$};\draw(1.5,1.6)circle(.16);\fill(1.5,1.6)circle(1.5pt);\node[right]at(1.8,1.6){$\vct B$朝外};\draw[->](.5,-.6)--(3,-.6)node[midway,below]{$\vct V$};}
图\ref{fig:m3fieldclasses}的方向构造还可从能流理解：正交场的坡印廷矢量沿$\vct E\times\vct B$。但$E=cB\ne0$时，两种尝试均要求$V=c$，没有可消去单个场的惯性系；平面电磁波属于这一类零不变量场。两个不变量都为零并不表示场本身为零。

当$\vct E\cdot\vct B\ne0$时，任一场在任何惯性系中都非零。在电、磁场平行的系中，两场的有向大小仍须同时满足两个不变量；确定场型后，再求相应的亚光速变换速度。

\section{从三力到轨迹}
对$\vct p=m\gamma_u\vct u$求导，可把三力写成
\[
 \vct f=m\gamma_u\vct a+m\gamma_u^3
 \frac{\vct u\cdot\vct a}{c^2}\vct u.
\]
平行于瞬时速度的投影给$f_\parallel=m\gamma_u^3a_\parallel$，垂直投影给$f_\perp=m\gamma_ua_\perp$。再反解为
\[
 \vct a=\frac1{m\gamma_u}
 \left[\vct f-\frac{\vct f\cdot\vct u}{c^2}\vct u\right].
\]
这个矢量式比引入方向依赖的“质量”更直接：静质量$m$是同一不变量，方向差异出自动量与速度的非线性关系。

\begin{example}\label{ex:m3transconstant}
粒子初始动量为$p_0\hat{\vct x}$，从$t=0$起受固定方向恒定三力$f\hat{\vct y}$。求速度和轨迹，说明为何它一般不是匀速圆周运动。
\Diagram{fig:m3transq}{固定方向力只在初始瞬间垂直速度}{\draw[axis](0,0)--(3.7,0)node[right]{$p_0$};\draw[axis](0,0)--(0,2.6)node[above]{$f$};\fill(0,0)circle(2pt);}
\end{example}
\begin{solution}
图\ref{fig:m3transq}中，三力固定沿$y$，故$\dd p_x/\dd t=0$、$\dd p_y/\dd t=f$，初值给$p_x=p_0$、$p_y=ft$。记$E_a=\sqrt{m^2c^4+p_0^2c^2}$，由质壳关系
\[
 E=\sqrt{E_a^2+f^2c^2t^2},\qquad
 u_x=\frac{c^2p_0}{E},\quad u_y=\frac{c^2ft}{E}.
\]
令$s=fct/E_a$，位移的两个定积分分别为
\begin{align*}
 x&=\frac{cp_0}f\int_0^s\frac{\dd q}{\sqrt{1+q^2}}
 =\frac{cp_0}f[\operatorname{arsinh}q]_0^s
 =\frac{cp_0}f\operatorname{arsinh}s,\\
 y&=\frac{E_a}f\int_0^s\frac{q\dd q}{\sqrt{1+q^2}}
 =\frac{E_a}f[\sqrt{1+q^2}]_0^s
 =\frac{E_a}f(\sqrt{1+s^2}-1).
\end{align*}
\Diagram{fig:m3transa}{动量沿直线增长，空间轨迹逐渐转向力的方向}{\draw[axis](0,0)--(4,0)node[right]{$x$};\draw[axis](0,0)--(0,4.1)node[above]{$y$};\draw[curve,domain=0:3.5,samples=80]plot({ln(\x+sqrt(1+\x*\x))},{sqrt(1+\x*\x)-1});\draw[->](1.8,2)--(1.95,2.6);}
当$p_0\ne0$时消去$s$，空间轨迹为
\[
 y=\frac{E_a}f\left[\cosh\left(\frac{fx}{cp_0}\right)-1\right].
\]
图\ref{fig:m3transa}中的曲线逐渐转向力的方向。由上式$E-E_a=fy$，可直接核对恒力做功；只在初始瞬间力与速度正交，随后$fu_y>0$，速率增大，因此不是匀速圆周运动。短时展开给$x\simeq c^2p_0t/E_a$、$y\simeq fc^2t^2/(2E_a)$，初始横向加速度为$f/(m\gamma_{u0})$。若$p_0=0$，粒子始终沿$y$直线运动，含$1/p_0$的消参式不适用，应使用原来的参数解。
\end{solution}

\begin{example}\label{ex:m3massgain}
质量$M$的静止物体完全吸收一个能量$E$的光子，不向外辐射。求吸收后物体静质量与速度；为什么不能把其静质量简单写成$M+E/c^2$？
\Diagram{fig:m3gainq}{吸收一个光子同时带入能量和动量}{\draw[->,thick](-2,.5)--(0,.5)node[midway,above]{$E$};\draw[curve](.3,0)rectangle(2,1);\node at(1.15,.5){$M$};}
\end{example}
\begin{solution}
图\ref{fig:m3gainq}中，入射光子动量为$E/c$，吸收后物体必须反冲。全系统总能量、动量分别为$E_T=Mc^2+E$、$P_T=E/c$，故
\[
 M_f^2c^4=E_T^2-c^2P_T^2=M^2c^4+2Mc^2E.
\]
令$\epsilon=E/(Mc^2)$，得到$M_f=M\sqrt{1+2\epsilon}$、$V=c\epsilon/(1+\epsilon)$。后式由$V=P_Tc^2/E_T$求出，不能把总动量除以末静质量当作速度。
\Diagram{fig:m3gaina}{光能一部分成为物体静能，一部分成为整体反冲动能}{\draw[curve](0,0)rectangle(5,1);\draw[curve](3.4,0)--(3.4,1);\node at(1.7,.5){$(M_f-M)c^2$};\node at(4.2,.5){$K_f$};\node[above]at(2.5,1){总输入$E$};}
图\ref{fig:m3gaina}的两部分分别为
\[
 \Delta U=Mc^2(\sqrt{1+2\epsilon}-1),\quad
 K_f=Mc^2(1+\epsilon-\sqrt{1+2\epsilon}).
\]
后一括号可有理化为$\epsilon^2/[1+\epsilon+\sqrt{1+2\epsilon}]$，因此正且在低能时为$\epsilon^2/2$。于是$\Delta U\simeq E-E^2/(2Mc^2)$、$K_f\simeq E^2/(2Mc^2)$，总和严格为$E$。

若错误地取$M_f=M+E/c^2$，便已把全部光能算入静能，却又保留非零末动量，末态总能量会超过初态。两束等能反向光子同时吸收时总动量为零，才可把全部输入能量计为静能。对吸收所致的变质量过程，$u\cdot F=c^2\dd M/\dd\tau$一般非零，与静质量不变点电荷的洛伦兹四力应区别。
\end{solution}

\section{电荷、电流与场的同一变换}
\begin{example}\label{ex:m3movingcharge}
点电荷$q$沿$x$以$V$匀速运动，在$t=0$经过原点。由静止库仑场求实验室时刻$t$、位置$(x,y,z)$的电磁场，描述纵向与横向方向的差别。
\Diagram{fig:m3chargeq}{相对电荷同时位置的矢量$\vct R$}{\draw[axis](-1,0)--(4.5,0)node[right]{$x$};\fill(1,0)circle(2pt)node[below]{$q,\ Vt$};\fill(3,1.8)circle(1.5pt)node[right]{$P$};\draw[->](1,0)--(3,1.8)node[midway,above]{$\vct R$};\draw[axis](1,.4)--(2,.4)node[above]{$V$};}
\end{example}
\begin{solution}
取图\ref{fig:m3chargeq}的$R_x=x-Vt$、$\vct R_\perp=(y,z)$。变到电荷静止系，场点相对电荷坐标为$(\gamma R_x,y,z)$，距离$r'=\sqrt{\gamma^2R_x^2+R_\perp^2}$，库仑场$\vct E'=q(\gamma R_x,y,z)/(4\pi\varepsilon_0r'^3)$、$\vct B'=0$。逆场变换保持平行分量，横向分量乘$\gamma$，所以
\[
 \vct E=\frac{q\gamma\vct R}{4\pi\varepsilon_0(\gamma^2R_x^2+R_\perp^2)^{3/2}},\quad
 \vct B=\frac{\vct V\times\vct E}{c^2}.
\]
\Diagram{fig:m3chargea}{等距球面上横向电场增强，纵向减弱}{\draw[axis](-2.4,0)--(2.7,0)node[right]{$V$};\draw[guide](0,0)circle(1.5);\draw[axis](0,1.5)--(0,2.7)node[above]{$\gamma E_C$};\draw[axis](1.5,0)--(2,0);\node[below]at(1.8,-.3){$E_C/\gamma^2$};\fill(0,0)circle(2pt);}
沿纵向$R_\perp=0$，$E=E_C/\gamma^2$；纯横向$R_x=0$，$E=\gamma E_C$，其中$E_C=q/(4\pi\varepsilon_0R^2)$。图\ref{fig:m3chargea}说明角分布被压向横向。场沿同时位置的$\vct R$指向并不意味着作用瞬时传播；这是匀速轨迹的完整延迟场化简结果，电荷加速时不能继续照用。
\end{solution}

\begin{example}\label{ex:m3wireforce}
实验室电中性无限直导线沿$x$，电流$I>0$。测试电荷$q$位于$y=r>0$并以$V\hat{\vct x}$运动。求实验室和电荷瞬时静止系中的横向力，用导线的变换电荷密度解释后一结果。
\Diagram{fig:m3wireforceq}{导线上方的磁场沿$+z$，正电荷受力向下}{\draw[axis](-2,0)--(3,0)node[right]{$I$};\fill(0,1.7)circle(2pt)node[left]{$q$};\draw[axis](0,1.7)--(1.6,1.7)node[above]{$V$};\draw(.5,.8)circle(.12);\fill(.5,.8)circle(.8pt);\node[right]at(.8,.8){$B_z$};}
\end{example}
\begin{solution}
图\ref{fig:m3wireforceq}中$B_z=\mu_0I/(2\pi r)$、$E_y=0$，实验室$f_y=-qVB_z$。变到电荷系，$E_y'=-\gamma VB_z$，电荷静止只受电力，故$f_y'=qE_y'=-\gamma qVB_z=\gamma f_y$，与横向三力变换一致。
\Diagram{fig:m3wireforcea}{电荷系导线带负线密度，产生向导线的电场}{\draw[curve](-2,0)--(3,0);\foreach\x in{-1.8,-1.1,-.4,.3,1,1.7,2.4}{\node at(\x,.25){$-$};}\fill(0,1.7)circle(2pt);\draw[axis](0,1.5)--(0,.5)node[right]{$E_y'<0$};}
由四电流变换，$\lambda'=-\gamma VI/c^2$。无限线电荷在图\ref{fig:m3wireforcea}处电场为
\[
 E_y'=\lambda'/(2\pi\varepsilon_0r)=-\gamma V\mu_0I/(2\pi r),
\]
与场变换完全相同。这个模型把实验室的纯磁力解释为电荷系的电力，但不意味着可以在所有空间点、所有电荷速度下选同一个系消去全部磁场。
\end{solution}

\begin{example}\label{ex:m3drift}
正交均匀场$\vct E=E\hat{\vct y}$、$\vct B=B\hat{\vct z}$满足$E<cB$。求可匀速直线通过而不偏转的带电粒子速度，并从消电场参考系解释。
\Diagram{fig:m3driftq}{速度选择器中电力与磁力相反}{\draw[axis](0,0)--(3.7,0)node[right]{$v_d$};\draw[axis](1,0)--(1,1.5)node[right]{$qE$};\draw[axis](1,0)--(1,-1.5)node[right]{$qv_dB$};\draw(2,1)circle(.15);\fill(2,1)circle(1pt);\node[right]at(2.3,1){$B$};}
\end{example}
\begin{solution}
图\ref{fig:m3driftq}中令一般速度为$(v_x,v_y,v_z)$。叉乘逐分量为$\vct v\times\vct B=(Bv_y,-Bv_x,0)$，零力条件给
\[
 Bv_y=0,\qquad E-Bv_x=0.
\]
因此全部无偏转速度构成一族：
\[
 \vct v=\frac EB\hat{\vct x}+v_z\hat{\vct z},\qquad
 |v_z|<\sqrt{c^2-\frac{E^2}{B^2}}.
\]
沿磁场的速度分量不产生磁力，因此保留为自由参数。只有额外规定$v_z=0$时，才得到唯一的二维选择速度$E/B$。
\Diagram{fig:m3drifta}{消电场系中的无偏转速度沿磁场；实验室还具有共同横向漂移}{\begin{scope}\draw[axis](0,0)--(2.6,0)node[right]{$x$};\draw[axis](0,-1.4)--(0,1.8)node[above]{$z$};\draw[->,thick](0,0)--(1.7,1.1)node[right]{$\vct v$};\draw[guide](0,1.1)--(1.7,1.1)--(1.7,0);\node[below]at(1.4,0){$E/B$};\node[left]at(0,.9){$v_z$};\end{scope}\begin{scope}[xshift=4cm]\draw[axis](0,0)--(2,0)node[right]{$x\prime$};\draw[axis](0,-1.4)--(0,1.8)node[above]{$z\prime$};\draw[->,thick](0,0)--(0,1.2)node[right]{$v_z\prime$};\draw[->](1.1,-.6)--(1.1,1)node[above]{$\vct B\prime$};\node[below]at(.7,-1){$\vct E\prime=0$};\end{scope}}
取消电场系沿$+x$以$V_d=E/B$运动，记$\gamma_d=(1-V_d^2/c^2)^{-1/2}$。场变换给
\[
 E'_y=\gamma_d(E-V_dB)=0,\qquad
 B'_z=\gamma_d(B-V_dE/c^2)=B/\gamma_d.
\]
对上述速度族，$v'_x=v'_y=0$，而
\[
 v'_z=\frac{v_z}{\gamma_d(1-V_d^2/c^2)}=\gamma_dv_z.
\]
图\ref{fig:m3drifta}中粒子在新系沿$\vct B'$直行，仍不受力；$|v'_z|<c$正好给出实验室的速度界限。若新系另有垂直于$\vct B'$的速度分量，粒子就会回旋，那是有偏转运动，已不属于题目所求的无偏转速度族。$E\geq cB$时不存在有质量粒子的这一亚光速速度族。
\end{solution}

\section{散射与衰变中的不变质量}
\begin{example}\label{ex:m3decay}
静质量$M$的粒子在自身静止系衰变为两个光子。实验室母粒子以速度$V$运动。求两光子实验室能量的范围；若源系两光子都垂直运动轴，求实验室夹角。
\Diagram{fig:m3decayq}{母粒子系两光子能量相等、方向相反}{\draw[->,thick](0,0)--(2,1.3)node[right]{$Mc^2/2$};\draw[->,thick](0,0)--(-2,-1.3)node[left]{$Mc^2/2$};\draw[axis](0,0)--(3,0)node[right]{$V$};}
\end{example}
\begin{solution}
图\ref{fig:m3decayq}静止系动量守恒给两光子能量均为$E_*=Mc^2/2$。若第一光子源系方向余弦为$\mu_*$，则$E_1=\gamma E_*(1+\beta\mu_*)$、$E_2=\gamma E_*(1-\beta\mu_*)$，两者之和为$\gamma Mc^2$。单个能量范围为$[\gamma E_*(1-\beta),\gamma E_*(1+\beta)]$。
\Diagram{fig:m3decaya}{源系横向光子在实验室形成前向对称双束}{\draw[axis](0,0)--(4,0)node[right]{$x$};\draw[->,thick](0,0)--(3,1.6);\draw[->,thick](0,0)--(3,-1.6);\node[right]at(3,1.6){$E_1=\gamma E_*$};\node[right]at(3,-1.6){$E_2=\gamma E_*$};}
若$\mu_*=0$，每个光子在实验室有$\cos\theta=\beta$，夹角为$2\arccos\beta$，如图\ref{fig:m3decaya}。再由两光子不变质量$2E_1E_2(1-\cos2\theta)=M^2c^4$核对：$1-\cos2\theta=2/\gamma^2$，正好恢复母粒子质量。
\end{solution}

\begin{example}\label{ex:m3nuclearpair}
光子撞击静止核，产生电子正电子对，核最终仍为原静质量$M$，电子质量$m$。忽略束缚能，求光子能量阈值，解释核反冲不能完全忽略的原因。
\Diagram{fig:m3nuclearq}{核提供反冲动量，末态总静质量至少$M+2m$}{\draw[->,thick](-2,0)--(0,0)node[midway,above]{$E$};\fill(.5,0)circle(4pt)node[below]{$M$};\draw[->](1,0)--(2,0);\foreach\x/\t in{2.7/M,3.7/{e^-},4.7/{e^+}}{\fill(\x,0)circle(2pt)node[below]{$\t$};}}
\end{example}
\begin{solution}
图\ref{fig:m3nuclearq}初态不变量$s=M^2c^4+2Mc^2E$。阈值时末态三粒子在质心系无相对动量，总质心能量$(M+2m)c^2$。因此
\[
 M^2c^4+2Mc^2E_{\rm th}=(M+2m)^2c^4,
\quad E_{\rm th}=2mc^2\left(1+\frac mM\right).
\]
\Diagram{fig:m3nucleara}{阈值在质心系三粒子同速，实验室仍须整体反冲}{\draw[curve](0,-.5)rectangle(4,.7);\foreach\x/\t in{.7/M,2/{e^-},3.3/{e^+}}{\fill(\x,0)circle(2pt)node[above]{$\t$};}\draw[axis](4.2,.1)--(5.5,.1)node[above]{$V_{\rm cm}$};}
大核质量极限才趋$2mc^2$。图\ref{fig:m3nucleara}中阈值末态在实验室具有共同质心速度，额外项提供反冲能量。若题设改为电子作靶且末态多一个电子，必须重新计总静质量，不能只把$M=m$代入后忽略相同粒子与末态计数。
\end{solution}

\begin{example}\label{ex:m3fourtensor}
在本书度规约定下，构造使$F^\mu=q\mathcal F^{\mu\nu}u_\nu$再现洛伦兹三力的反对称场张量，并说明反对称性怎样保证静质量不变。
\Diagram{fig:m3tensorq}{同一洛伦兹力的时间分量是功率，空间分量是三力}{\draw[axis](0,0)--(3.2,0)node[right]{$\vct u$};\draw[axis](0,0)--(1.3,2.1)node[above]{$\vct E$};\draw(2,1)circle(.15);\fill(2,1)circle(1pt);\node[right]at(2.3,1){$\vct B$};}
\end{example}
\begin{solution}
按$(ct,x,y,z)$排序、度规取$(+,-,-,-)$，图\ref{fig:m3tensorq}所用场张量为
\[
 \mathcal F^{\mu\nu}=\begin{pmatrix}
 0&-E_x/c&-E_y/c&-E_z/c\\
 E_x/c&0&-B_z&B_y\\
 E_y/c&B_z&0&-B_x\\
 E_z/c&-B_y&B_x&0
 \end{pmatrix}.
\]
粒子四速度降指标后为$u_\nu=\gamma_u(c,-u_x,-u_y,-u_z)$。矩阵第0行给$F^0=q\gamma_u\vct E\cdot\vct u/c$；第$x$行给
\[
 F^x=q\gamma_u(E_x+u_yB_z-u_zB_y).
\]
其余两行循环置换指标，得到$\vct F=q\gamma_u(\vct E+\vct u\times\vct B)$。再除以$\gamma_u$便还原洛伦兹三力。
\Diagram{fig:m3tensora}{反对称张量与对称速度乘积缩并为零}{\draw[curve](0,0)rectangle(2.6,1.1);\node at(1.3,.55){$\mathcal F^{\mu\nu}=-\mathcal F^{\nu\mu}$};\draw[->](2.8,.55)--(3.7,.55);\node[right]at(3.8,.55){$u_\mu F^\mu=0$};}
图\ref{fig:m3tensora}中的正交关系来自反对称性。设$Q=u_\mu\mathcal F^{\mu\nu}u_\nu$，交换两个求和哑指标得
\[
 Q=u_\nu\mathcal F^{\nu\mu}u_\mu
 =-u_\mu\mathcal F^{\mu\nu}u_\nu=-Q,
\]
故$Q=0$。为说明它与静质量的关系，暂不预设$m$恒定，将$p^\mu=mu^\mu$代入：
\[
 u_\mu\frac{\dd p^\mu}{\dd\tau}
 =c^2\frac{\dd m}{\dd\tau}
 +\frac m2\frac{\dd(u_\mu u^\mu)}{\dd\tau}
 =c^2\frac{\dd m}{\dd\tau}.
\]
左端由洛伦兹四力为零，所以$\dd m/\dd\tau=0$。电磁场可改变点电荷的运动能量和动量，但在此无内部激发的模型中不改变其静质量。
\end{solution}

\section{弹性碰撞的能量转移与交叉场轨迹}
\subsection{固定靶碰撞中的最大反冲能量}
设入射粒子静质量$M$、实验室能量$E=\gamma Mc^2$、动量$p$，静止靶粒子质量$m$，碰撞前后两粒子的静质量不变。用$s$记总四动量平方乘$c^2$，则
\[
 s=M^2c^4+m^2c^4+2mc^2E,
 \qquad V_* =\frac{pc^2}{E+mc^2},\quad
 \gamma_* =\frac{E+mc^2}{\sqrt s}.
\]
质心系中两粒子动量大小为$p_*$。把两体公式的根号因式展开，利用$E^2-p^2c^2=M^2c^4$，得到
\[
 p_*^2c^2=\frac{m^2c^4(E^2-M^2c^4)}s,
 \qquad p_* =\frac{mc^2p}{\sqrt s}.
\]
靶粒子在质心系中的能量固定，实验室能量随其质心系出射角$\vartheta_*$变化：$E_m=\gamma_*(E_{m*}+V_*p_*\cos\vartheta_*)$。初态靶在该系向后运动，实验室能量为$mc^2=\gamma_*(E_{m*}-V_*p_*)$。两式相减，直接得到反冲动能
\[
 T=\gamma_*V_*p_*(1+\cos\vartheta_*).
\]
最大值对应靶沿质心系正前方射出，而不是任意把实验室两粒子都假定为向前且等速。代入各参数有
\[
 T_{\max}=\frac{2m c^2(\gamma^2-1)}
 {1+2\gamma m/M+(m/M)^2}.
\]
这是纯运动学上界；相互作用的微分截面可能使达到上界的事件很少，但不改变该上界的位置。

\begin{example}重粒子对轻靶的最大能量转移
入射粒子质量为$M=100m$，洛伦兹因子为2，靶粒子初始静止。求最大反冲能量占入射动能的比例；再检查同质量碰撞和非相对论极限。
\ExpansionDiagram{fig:dep-m3-transfer-q}{实验室固定靶与质心系反冲端点}{
\draw[->,thick](0,0)--(2,0)node[midway,above]{$p$};\fill(2.4,0)circle(2pt)node[below]{$M$};\fill(4.5,0)circle(2pt)node[below]{$m$静止};\draw[dashed](5.6,-.6)--(5.6,1.2);\draw[->](6.2,0)--(8,0)node[midway,above]{$p_*$};\node[below]at(7,-.4){质心系靶正向};}
\end{example}
\begin{solution}
由图\ref{fig:dep-m3-transfer-q}的正向质心系端点，
\[
 T_{\max}=\frac{6mc^2}{1+0.04+0.0001}
 \simeq5.7687mc^2.
\]
入射动能$K=(\gamma-1)Mc^2=100mc^2$，最大转移比例约为$5.77\%$。尽管入射动能很大，质量相差悬殊仍使反冲受限；能量和动量必须由同一组末态同时满足。
\ExpansionDiagram{fig:dep-m3-transfer-a}{反冲能量对质心系角度的依赖}{
\draw[->](0,0)--(5,0)node[right]{$\cos\vartheta_*$};\draw[->](0,0)--(0,2.6)node[above]{$T/T_{\max}$};\draw[thick](.5,0)--(4.5,2);\node[below]at(.5,0){$-1$};\node[below]at(2.5,0){$0$};\node[below]at(4.5,0){1};\node[left]at(0,2){1};}
图\ref{fig:dep-m3-transfer-a}的直线关系为$T/T_{\max}=(1+\cos\vartheta_*)/2$，只描述允许能量与角度的对应，不表示各角度等概率。

若$M=m$，分母为$2(1+\gamma)$，所以$T_{\max}=mc^2(\gamma-1)=K$，同质量的一维弹性碰撞可以交换运动状态。若入射速度很小，$\gamma^2-1\simeq v^2/c^2$，得到$T_{\max}\simeq2mM^2v^2/(M+m)^2$；再除以$Mv^2/2$，恢复经典比例$4mM/(M+m)^2$。这两个极限分别检查相同质量和低速度的退化情形。
\end{solution}

\subsection{消去电场后仍须变换整条世界线}
磁占优的正交场可以找到纯磁参考系，但这不意味着实验室中粒子必作匀速直线运动。只有粒子在那个参考系恰好静止时，才对应无偏转速度。一般初态仍有回旋运动，变回实验室后既有漂移又有速率变化。

\begin{example}由静止释放的粒子在正交场中运动
实验室有$\vct E=E\hat{\vct y}$、$\vct B=B\hat{\vct z}$，其中$0<E<cB$。质量$m$、电量$q>0$的粒子在原点由静止释放。忽略辐射，求参数轨迹、平均漂移速度和最大动能。
\ExpansionDiagram{fig:dep-m3-cross-q}{磁场朝纸外，电场向上，无偏转漂移方向向右}{
\draw[->](0,0)--(5,0)node[right]{$x$};\draw[->](0,0)--(0,2.7)node[above]{$y$};\draw[->,thick](1,1)--(1,2)node[above]{$E$};\draw(3,1.5)circle(.18);\fill(3,1.5)circle(1.4pt)node[right=7pt]{$B$};\fill(0,0)circle(2pt)node[below left]{$q$};\draw[->](1,-.6)--(3,-.6)node[midway,below]{$V=E/B$};}
\end{example}
\begin{solution}
按图\ref{fig:dep-m3-cross-q}取$S'$相对实验室以$V=E/B$沿$+x$运动，记$\beta=V/c$、$\Gamma=(1-\beta^2)^{-1/2}$。场变换给$E'=0$、$B'=B/\Gamma$。粒子初始在$S'$的速度为$-V\hat{\vct x}$，所以其恒定洛伦兹因子是$\Gamma$，回旋角频率为
\[
 \Omega=\frac{qB'}{\Gamma m}=\frac{qB}{\Gamma^2m}.
\]
洛伦兹力分量为$\dot v'_x=\Omega v'_y$、$\dot v'_y=-\Omega v'_x$。将第一式再求导并代入第二式，得$\ddot v'_x+\Omega^2v'_x=0$，通解为$A\cos\Omega t'+B\sin\Omega t'$。初值$v'_x(0)=-V$、$\dot v'_x(0)=\Omega v'_y(0)=0$分别给$A=-V,B=0$；再由$v'_y=\dot v'_x/\Omega$得$v'_y=V\sin\Omega t'$。令$s=\Omega t'$，从原点积分：
\begin{align*}
 x'&=\int_0^{t'}-V\cos(\Omega q)\dd q
 =\left[-\frac V\Omega\sin(\Omega q)\right]_0^{t'},\\
 y'&=\int_0^{t'}V\sin(\Omega q)\dd q
 =\left[-\frac V\Omega\cos(\Omega q)\right]_0^{t'}.
\end{align*}
代入两端后，
\[
 x'=-\frac V\Omega\sin s,\quad
 y'=\frac V\Omega(1-\cos s),\quad t'=s/\Omega.
\]
逆变换必须同时作用于时间和位置，得到
\[
 x=\frac{\Gamma V}{\Omega}(s-\sin s),\qquad
 y=\frac V\Omega(1-\cos s),\qquad
 t=\frac\Gamma\Omega(s-\beta^2\sin s).
\]
\ExpansionDiagram{fig:dep-m3-cross-a}{实验室轨迹的无量纲摆线；每个尖点处粒子瞬时静止}{
\draw[->](0,0)--(7,0)node[right]{$x\Omega/(\Gamma V)$};\draw[->](0,0)--(0,2.7)node[above]{$y\Omega/V$};\draw[thick,domain=0:6.283185,samples=180]plot({\x-sin(\x r)},{1-cos(\x r)});\fill(0,0)circle(1.4pt);\fill(6.283185,0)circle(1.4pt);\node[below]at(6.283185,0){$2\pi$};}
图\ref{fig:dep-m3-cross-a}中横纵轴使用不同长度单位，轨迹形状对应上述参数式。因为$\dd t/\dd s=\Gamma(1-\beta^2\cos s)/\Omega>0$，参数没有倒转时间。在一周$s:0\to2\pi$内，$\Delta x=2\pi\Gamma V/\Omega$、$\Delta t=2\pi\Gamma/\Omega$，所以平均漂移为$V=E/B$。

磁力不做功，实验室动能等于$qEy$，于是
\[
 K(s)=\frac{qEV}{\Omega}(1-\cos s)
 =\Gamma^2mV^2(1-\cos s),\quad
 K_{\max}=2\Gamma^2mV^2.
\]
也可变换四动量得粒子能量$\Gamma^2mc^2(1-\beta^2\cos s)$，减去$mc^2$得到同一结果。一周后高度回零，动能也回零，但位置已漂移；这与“在纯磁系能量恒定”相容，两个系的能量定义不同。
\end{solution}

\section{两体末态的能谱与连续光压推进}
\subsection{各向同性衰变怎样形成均匀能谱}
母粒子静质量$M$，在自身静止系衰变为两个同质量$m$的粒子。每个子粒子的$E_*=Mc^2/2$、$p_*c=\sqrt{M^2/4-m^2}\,c^2$。若衰变在该系各向同性，单个已标记子粒子的$\mu_*=\cos\theta_*$在$[-1,1]$上均匀。

母粒子在实验室沿$x$以$\beta c$运动，洛伦兹因子为$\Gamma$。子粒子能量为
\[
 E=\Gamma(E_*+\beta cp_*\mu_*).
\]
在$\beta p_*>0$时，这是$\mu_*$的严格递增线性函数，由$\dd P=\dd\mu_*/2$及$\dd E=\Gamma\beta cp_*\dd\mu_*$，全角接收的能量概率密度为
\[
 \frac{\dd P}{\dd E}=\frac1{2\Gamma\beta cp_*},\qquad
 E_-\leq E\leq E_+,\quad
 E_\pm=\Gamma(E_*\pm\beta cp_*).
\]
若$\beta=0$或$p_*=0$，两端点合并为$E=\Gamma E_*$，成为单能分布，不能再使用上式的零分母密度。实验室角分布另含方向变换的雅可比；有限角域探测须先反解源系接受范围，再截取能谱。

\begin{example}由能谱端点恢复子粒子质量
已知母粒子实验室洛伦兹因子为2，两体等质量衰变的单子粒子全角能谱范围为$4\mathcal E\leq E\leq16\mathcal E$，$\mathcal E$为能量单位。假定母粒子静止系中衰变方向各向同性。求母、子静质量，求$E>12\mathcal E$的概率，并说明同一次衰变中两子粒子能量的相关性。
\ExpansionDiagram{fig:dep-m3-spectrum-q}{全角接收的两体末态形成有两个运动学端点的能谱}{
\draw[->](0,0)--(7,0)node[right]{$E/\mathcal E$};\draw[->](0,0)--(0,2.5)node[above]{概率密度};\draw[thick](1.5,0)--(1.5,1.5)--(6,1.5)--(6,0);\node[below]at(1.5,0){4};\node[below]at(6,0){16};\node[above]at(3.8,1.5){均匀};}
\end{example}
\begin{solution}
图\ref{fig:dep-m3-spectrum-q}的中点为$E_c=10\mathcal E=\Gamma Mc^2/2$，所以$Mc^2=10\mathcal E$。半宽$\Delta=6\mathcal E=\Gamma\beta cp_*$，而$\Gamma\beta=\sqrt3$，得$cp_*=2\sqrt3\mathcal E$。由$E_*^2=m^2c^4+p_*^2c^2$，
\[
 m^2c^4=(5\mathcal E)^2-(2\sqrt3\mathcal E)^2
 =13\mathcal E^2,\qquad mc^2=\sqrt{13}\mathcal E.
\]
矩形分布中，高于$12\mathcal E$的区间占长度$4\mathcal E$，总宽为$12\mathcal E$，所以概率$1/3$。
\ExpansionDiagram{fig:dep-m3-spectrum-a}{同一次衰变的能量对被限制在总能量直线上}{
\draw[->](0,0)--(5,0)node[right]{$E_1$};\draw[->](0,0)--(0,4.5)node[above]{$E_2$};\draw[thick](1,4)--(4,1);\fill(1,4)circle(1.4pt);\fill(4,1)circle(1.4pt);\node[right]at(3.3,3.4){$E_1+E_2=20\mathcal E$};}
图\ref{fig:dep-m3-spectrum-a}中$E_2=20\mathcal E-E_1$，两者完全反相关。令$\xi=E_1-E_c$，则$\dd E_1=\dd\xi$，积分限为$-\Delta,\Delta$，密度为$1/(2\Delta)$。对称性给$\langle\xi\rangle=0$，方差为
\[
 \operatorname{Var}E_1=\frac1{2\Delta}\int_{-\Delta}^{\Delta}\xi^2\dd\xi
 =\left[\frac{\xi^3}{6\Delta}\right]_{-\Delta}^{\Delta}
 =\frac{\Delta^2}{3}=12\mathcal E^2.
\]
又有$E_2-E_c=-(E_1-E_c)$，故协方差为$-\langle\xi^2\rangle=-12\mathcal E^2$。若把两个均匀谱独立抽样，会产生不守恒的能量对。若$\Gamma$也未知，两个端点只约束$\Gamma M$及$\beta\sqrt{1-4m^2/M^2}$，通常不足以唯一确定三个未知量。
\end{solution}

\subsection{移动镜面的光压还要计算追上镜面的光子率}
平面镜沿光束方向后退，实验室固定截面上的入射功率为$P$。若镜速为$\beta c$，单位实验室时间追上镜面的入射能量是$P(1-\beta)$，反射单光子能量比为$(1-\beta)/(1+\beta)$。每个光子反向弹回，镜获得的动量是入射与反射光子动量模之和，因此
\[
 F=\frac{P(1-\beta)}c\left(1+\frac{1-\beta}{1+\beta}\right)
 =\frac{2P}{c}\frac{1-\beta}{1+\beta}.
\]
这里忽略单光子的有限反冲，采用大量弱光子的连续光压极限。若遗漏追赶率$1-\beta$，会高估高速镜所受的力。

\begin{example}恒定入射功率推动理想光帆
质量$M$的理想反射光帆初始静止，实验室光源提供固定截面功率$P$的连续准直光束，光帆全部截获该束并沿束方向运动，静质量不变。忽略衍射、束宽变化及辐射吸收。求速度和位置随时间的参数式，并核对光子能量损失与光帆做功。
\ExpansionDiagram{fig:dep-m3-sail-q}{光帆后退使光子到达率降低，反射光同时红移}{
\draw[->,thick](0,.7)--(4,.7)node[midway,above]{$P$};\draw[very thick](4,-.5)--(4,1.7);\draw[->](4,1.9)--(5.5,1.9)node[right]{$\beta c$};\draw[->](4,-.1)--(1,-.1)node[midway,below]{红移反射};\node[below]at(4,-.6){光帆$M$};}
\end{example}
\begin{solution}
图\ref{fig:dep-m3-sail-q}中用快度$\chi$表示$\beta=\tanh\chi$，有$(1-\beta)/(1+\beta)=\ee^{-2\chi}$。动量$Mc\sinh\chi$的导数满足
\[
 Mc\cosh\chi\frac{\dd\chi}{\dd t}=\frac{2P}c\ee^{-2\chi}.
\]
因此$\dd t/\dd\chi=Mc^2\ee^{2\chi}\cosh\chi/(2P)$。展开双曲函数，并用初值$t(0)=0$，
\[
 t=\frac{Mc^2}{4P}\int_0^\chi(\ee^{3q}+\ee^q)\dd q
 =\frac{Mc^2}{4P}\left[\frac{\ee^{3q}}3+\ee^q\right]_0^\chi.
\]
令$u=\ee^\chi\geq1$，得
\[
 t=\frac{Mc^2}{12P}(u^3+3u-4),\qquad
 \beta=\frac{u^2-1}{u^2+1}.
\]
位置由$\dd x/\dd\chi=\beta c\dd t/\dd\chi$积分。用$\beta=\tanh\chi$及$x(0)=0$，
\[
 x=\frac{Mc^3}{4P}\int_0^\chi(\ee^{3q}-\ee^q)\dd q
 =\frac{Mc^3}{4P}\left[\frac{\ee^{3q}}3-\ee^q\right]_0^\chi,
\]
得到
\[
 x=\frac{Mc^3}{12P}(u^3-3u+2).
\]
\ExpansionDiagram{fig:dep-m3-sail-a}{连续推进的速度单调趋近光速，所需时间随快度迅速增长}{
\draw[->](0,0)--(5,0)node[right]{$u=\ee^\chi$};\draw[->](0,0)--(0,2.5)node[above]{$v/c$};\draw[thick,domain=1:4.5,samples=80]plot({\x-1},{2*(\x^2-1)/(\x^2+1)});\draw[dashed](0,2)--(4.5,2);\node[left]at(0,2){1};\node[below]at(0,0){1};}
图\ref{fig:dep-m3-sail-a}的$u=1$为初始时刻。小时间下$F\simeq2P/c$，恢复初始加速度$2P/(Mc)$；高速时$t$按$u^3$增长，而$1-\beta\simeq2/u^2$，始终不能在有限时间达到$c$。

每单位实验室时间反射离开的光能为$P_{\rm refl}=P(1-\beta)^2/(1+\beta)$。真正撞到光帆的入射能量率为$P_{\rm hit}=P(1-\beta)$，两者之差
\[
 P_{\rm hit}-P_{\rm refl}=\frac{2P\beta(1-\beta)}{1+\beta}=Fv,
\]
与光帆动能增长率一致。固定截面功率$P$和镜面实际截获功率不同；未追上光帆的光还在传播途中，不应直接算作已交给光帆的能量。
\end{solution}

\section{习题与解答}
\begin{exercise}\label{ex:m3electricdominant}
正交均匀场$\vct E=E\hat{\vct y}$、$\vct B=B\hat{\vct z}$满足$E>cB$。带正电粒子初始实验室速度为$V\hat{\vct x}$，其中$V=c^2B/E$。求其轨迹，忽略辐射和边界。
\Diagram{fig:m3edominantq}{电占优场可变为纯电场，初速度选为相应推动速度}{\draw[axis](0,0)--(3.7,0)node[right]{$V$};\draw[axis](0,0)--(0,2.6)node[above]{$E$};\draw(1.8,1)circle(.15);\fill(1.8,1)circle(1pt);\node[right]at(2.1,1){$B$};}
\end{exercise}
\begin{solution}
图\ref{fig:m3edominantq}沿$V$推动后$B'=0$，$E'=\sqrt{E^2-c^2B^2}$，粒子初始静止。设$\alpha=qE'/m$，沿用习题\ref{ex:m1fourwork}已逐项积分的恒力解，并将加速轴改为$y'$，得到
\[
 ct'=\frac{c^2}\alpha\sinh(\alpha\tau/c),\quad
 y'=\frac{c^2}\alpha[\cosh(\alpha\tau/c)-1],\quad x'=0.
\]
逆推动给$t=\gamma t'$、$x=Vt$、$y=y'$，消去固有时得到
\[
 y(t)=\frac{c^2}\alpha\left[\sqrt{1+(\alpha t/(\gamma c))^2}-1\right].
\]
\Diagram{fig:m3edominanta}{轨迹按$X=\alpha x/(\gamma Vc)$、$Y=\alpha y/c^2$归一化，虚线为渐近线}{\draw[axis](0,0)--(4.5,0)node[right]{$X$};\draw[axis](0,0)--(0,3.5)node[above]{$Y$};\draw[curve,domain=0:4,samples=80]plot(\x,{sqrt(1+\x*\x)-1});\draw[guide](1,0)--(4.2,3.2);\node[right]at(3.7,2.4){$Y=X-1$};}
图\ref{fig:m3edominanta}的长时轨迹趋于直线斜向传播。横向速度为$v_y=(c/\gamma)\tanh(\alpha\tau/c)$，所以$v^2=V^2+c^2\tanh^2(\alpha\tau/c)/\gamma^2<c^2$。$v_x$恒定不代表$x$向动量恒定，因为总$\gamma_v$仍在增加；磁力的$x$分量正好提供这部分动量变化。
\end{solution}

\begin{exercise}\label{ex:m3decaycone}
静质量$M>2m$的粒子衰变为两个相同质量$m$的粒子，母粒子实验室速度为$V$。求子粒子在质心系的速率$u_*$；当$V>u_*$时，求所有实验室出射方向被限制的最大半角。
\Diagram{fig:m3decayconeq}{母粒子系两子粒子动量等大反向}{\draw[axis](-2.5,0)--(3.5,0)node[right]{$V$};\draw[->,thick](0,0)--(1.4,1.8)node[right]{$m$};\draw[->,thick](0,0)--(-1.4,-1.8)node[left]{$m$};}
\end{exercise}
\begin{solution}
图\ref{fig:m3decayconeq}中，母粒子静止系总动量为零，两子粒子能量相同，故$E_*=Mc^2/2$。质壳关系给
\[
 p_*c=\sqrt{M^2c^4/4-m^2c^4},\qquad
 u_*=\frac{p_*c^2}{E_*}=c\sqrt{1-4m^2/M^2}.
\]
令$z=\cos\theta_*$，母粒子沿$x$以$V$运动。速度变换的共同分母为$1+Vu_*z/c^2>0$，相除后
\[
 \tan\theta=\frac{u_*\sqrt{1-z^2}}{\gamma(V+u_*z)}.
\]
在$V>u_*$时分母始终正，所有非零速度都向前。
\Diagram{fig:m3decayconea}{亚光速子粒子可全部集中于一个前向圆锥}{\draw[axis](0,0)--(4.4,0)node[right]{$V$};\draw[curve](0,0)--(3,1.4);\draw[curve](0,0)--(3,-1.4);\draw[guide](3,-1.4)--(3,1.4);\node[above]at(2.2,1.2){$\theta_{\max}$};}
对$-1<z<1$求对数导数，最大角满足
\[
 -\frac z{1-z^2}-\frac{u_*}{V+u_*z}=0
 \quad\Longrightarrow\quad Vz+u_*=0.
\]
取$z=-u_*/V$，则$\sqrt{1-z^2}=\sqrt{V^2-u_*^2}/V$、$V+u_*z=(V^2-u_*^2)/V$，代入得到图\ref{fig:m3decayconea}的
\[
 \tan\theta_{\max}=\frac{u_*}{\gamma\sqrt{V^2-u_*^2}}.
\]
两端$z=\pm1$均有$\theta=0$，内点驻值因而为全局最大。速度相等时则出现退化：$V=u_*$时，正后方发射的子粒子在实验室静止，方向未定义；其余均向前，出射角可趋近但达不到$\pi/2$。$V<u_*$时取$z=-1$便有负的纵向速度，出现后向出射，且可有横向出射。光子$u_*=c$不可能满足$V>u_*$，因而不能用这个有限前向圆锥代替全角光行差。
\end{solution}

\begin{exercise}\label{ex:m3fizeau}
透明介质在自身静止系的非色散折射率为$n$，光沿介质运动方向传播，介质相对实验室速度$V$。求实验室相速度，并展开至$V/c$一阶。
\Diagram{fig:m3fizeauq}{介质系光速为$c/n$，介质整体沿同向运动}{\draw[curve](0,-.6)rectangle(4,1);\draw[axis](.4,.2)--(2.8,.2)node[above]{$c/n$};\draw[axis](4.2,.2)--(5.5,.2)node[above]{$V$};}
\end{exercise}
\begin{solution}
图\ref{fig:m3fizeauq}的波面速度沿$x$，由相位变换或速度合成都有$v_{\rm ph}=(c/n+V)/(1+V/(nc))$。低速把分母倒数展开为$1-V/(nc)+O(V^2/c^2)$，得$v_{\rm ph}=c/n+V(1-1/n^2)+O(V^2/c)$。
\Diagram{fig:m3fizeaua}{拖曳系数在$n=1$时为零}{\draw[axis](0,0)--(4.5,0)node[right]{$n$};\draw[axis](0,0)--(0,2.8)node[above]{$1-1/n^2$};\draw[curve,domain=1:4,samples=70]plot(\x,{2.4*(1-1/(\x*\x))});\node[below]at(1,0){$1$};}
图\ref{fig:m3fizeaua}在$n=1$时恢复真空光速$c$。本题采用非色散近似；真实色散介质中，变换也改变频率，应在相应介质静止频率处取$n$，一阶展开还会出现色散修正。
\end{solution}

\begin{exercise}\label{ex:m3currentrest}
某区域总电荷密度为零、电流密度$\vct j=j\hat{\vct x}\ne0$。是否存在使总电流为零的惯性系？解释为什么四电流不总能写为某个实固有密度乘四速度。
\Diagram{fig:m3currentq}{等量正负电荷相反运动，密度相消而电流相加}{\foreach\x in{0,1,2,3}{\node at(\x,.5){$+$};\node at(\x,-.5){$-$};}\draw[axis](.5,1)--(2,1);\draw[axis](2,-1)--(.5,-1);}
\end{exercise}
\begin{solution}
图\ref{fig:m3currentq}中，总四电流为$J^\mu=(0,j,0,0)$，不变量$J^2=-j^2<0$。假设某惯性系有$\vct j'=0$，则$J'^2=c^2\rho'^2\geq0$，与不变量符号矛盾，所以任何方向的惯性变换都不能消去总电流。

若仅沿$x$变换，可逐项看到$\rho'=-\gamma Vj/c^2$、$j'_x=\gamma j$。电中性不保持，但电流仍不为零。
\Diagram{fig:m3currenta}{类空四电流无法通过亚光速变换转到纯时间方向}{\draw[axis](-2,0)--(2.6,0)node[right]{$j_x$};\draw[axis](0,-1.8)--(0,2.4)node[above]{$c\rho$};\draw[guide](-1.8,-1.8)--(2,2);\draw[guide](-1.8,1.8)--(1.8,-1.8);\draw[->,thick](0,0)--(2,0);}
图\ref{fig:m3currenta}的四电流位于光锥外。单一带电成分共同运动时，$J^\mu=\rho_0U^\mu$且$J^2=\rho_0^2c^2\geq0$；中性电流却可能由正负两组四流相加。例如实验室密度为$+\rho$的正电荷静止，密度为$-\rho$的负电荷以$-u\hat{\vct x}$运动，总电荷为零、总电流为$+\rho u\hat{\vct x}$。各组分别有自己的静止系，总和没有共同电流静止系。

更一般地，只有$J^2>0$的非零四电流才可变到纯时间方向；$J^2<0$可变到纯空间方向，$J^2=0$且非零时两种静止化都不可行。总四电流按带符号电荷加权，与物质总能量动量不同；其类空性质不要求任何载体超光速。
\end{solution}

\begin{exercise}\label{ex:m3obliquewave}
平面波实验室传播方向与$x$成60度，观察系以$0.6c$沿$x$运动。求频率比和新方向角，核对横向波矢不变。
\Diagram{fig:m3obliquewaveq}{斜向波矢的纵横分解}{\draw[axis](0,0)--(3.8,0)node[right]{$x$};\draw[axis](0,0)--(1.5,2.598)node[right]{$\vct k$};\draw[guide](1.5,0)--(1.5,2.598);\node at(.9,.4){60度};}
\end{exercise}
\begin{solution}
图\ref{fig:m3obliquewaveq}有$\gamma=1.25$、$\cos\theta=.5$，频率比$\omega'/\omega=1.25(1-.6\times.5)=.875$。方向余弦$\cos\theta'=(.5-.6)/(1-.3)=-1/7$，所以$\theta'\simeq98.21^\circ$，略向后。
\Diagram{fig:m3obliquewavea}{两系采用相同波矢比例，横向分量的高度相同}{
\draw[axis](-1.8,0)--(2.2,0)node[right]{$k'_x$};\draw[axis](0,0)--(0,3.2)node[above]{$k'_y$};
\draw[->,thick](0,0)--(-.375,2.598)node[above left]{$\vct k'$};
\draw[guide](-.375,0)--(-.375,2.598)--(1.5,2.598)--(1.5,0);
\draw[->,dashed](0,0)--(1.5,2.598)node[above right]{$\vct k$};
\node[right]at(1.6,1.7){$k'_y=k_y$};}
新系$\sin\theta'=\sqrt{48}/7=4\sqrt3/7$，因此$k'_y=(\omega'/c)\sin\theta'=(7/8)(\omega/c)(4\sqrt3/7)=\sqrt3\omega/(2c)=k_y$。图\ref{fig:m3obliquewavea}的角度必须用负余弦判断象限，不能把正弦反求成81.79度。
\end{solution}

\begin{exercise}\label{ex:m3recoilphoton}
静止质量$M$的理想镜被能量$E$的单个光子正向撞击，光子反向弹回，镜的静质量不变。求反射光能量和镜反冲动能。
\Diagram{fig:m3recoilq}{有限质量镜不能视为预先给定速度的无限重反射面}{\draw[very thick](2,-1)--(2,1.2);\draw[->](0,.5)--(1.9,.5)node[midway,above]{$E$};\draw[->](1.9,-.3)--(0,-.3)node[midway,below]{$E_r$};\node[right]at(2.3,0){$M$};}
\end{exercise}
\begin{solution}
图\ref{fig:m3recoilq}镜末能量$Mc^2+E-E_r$，末动量$(E+E_r)/c$。代入镜的质壳关系，
\[
 (Mc^2+E-E_r)^2-(E+E_r)^2=M^2c^4.
\]
展开得$2Mc^2(E-E_r)=4EE_r$，所以$E_r=E/[1+2E/(Mc^2)]$。
\Diagram{fig:m3recoila}{镜的反冲动能等于光子损失能量}{\draw[axis](0,0)--(4.3,0)node[right]{$E/(Mc^2)$};\draw[axis](0,0)--(0,2.6)node[above]{$E_r/E$};\draw[curve,domain=0:3.5,samples=70]plot(\x,{2.2/(1+2*\x)});}
反冲动能$K_M=E-E_r=2E^2/(Mc^2+2E)$。图\ref{fig:m3recoila}在$M\to\infty$趋无频移；有限质量时必红移。这里不能先选某个“镜速”代入两次多普勒式，因为速度在单次碰撞中改变，需用全程守恒确定。
\end{solution}

\begin{exercise}\label{ex:m3invariantfields}
在某系有$\vct E=E_0\hat{\vct y}$、$\vct B=2E_0\hat{\vct z}/c$。求消电场的速度、剩余磁场和该系中粒子静止时对应实验室的速度。
\Diagram{fig:m3invariantq}{磁占优正交场}{\draw[axis](0,0)--(3.5,0)node[right]{$x$};\draw[axis](0,0)--(0,2.1)node[above]{$E_0$};\draw(1.6,1)circle(.15);\fill(1.6,1)circle(1pt);\node[right]at(1.9,1){$2E_0/c$};}
\end{exercise}
\begin{solution}
图\ref{fig:m3invariantq}给出$E=E_0$、$B=2E_0/c$，且$\vct E\perp\vct B$，磁占优。沿$+x$变换，有$E_y'=\gamma(E_0-V\,2E_0/c)$。令其为零，得到$V=c/2$、$\gamma=2/\sqrt3$。再求剩余磁场，
\[
 B_z'=\frac2{\sqrt3}\left(\frac{2E_0}c-\frac{E_0}{2c}\right)
 =\frac{\sqrt3E_0}c.
\]
也可由$B'^2=B^2-E^2/c^2=3E_0^2/c^2$独立核对其大小，方向由具体变换保留为$+z$。
\Diagram{fig:m3invarianta}{消场系中静止对应实验室匀速漂移}{\fill(0,0)circle(2pt);\draw[axis](0,.5)--(2.5,.5)node[right]{$V=c/2$};\node[below]at(1,-.4){$\vct E+\vct V\times\vct B=0$};}
图\ref{fig:m3invarianta}中的粒子在新系静止，逆速度变换给实验室$\vct v=(c/2)\hat{\vct x}$。其磁力方向为$\hat{\vct x}\times\hat{\vct z}=-\hat{\vct y}$，大小$q(c/2)(2E_0/c)=qE_0$，正好与电力相抵。

若只要求新系粒子沿$z'$运动而非静止，则逆变换仍有$v_x=c/2$，还允许$v_z=v_z'/\gamma$，范围$|v_z|<\sqrt3c/2$。这进一步说明消电场不等于消去所有粒子的力：沿$\vct B'$运动仍零力，垂直$\vct B'$运动却受磁力并回旋。参考系选择由场决定，具体粒子是否无偏转还要代入其速度。
\end{solution}

\begin{exercise}\label{ex:m3properpower}
粒子静质量不变，某实验室三力始终平行瞬时速度。证明四力模等于三力模，并说明这并不意味着四力时间分量为零。
\Diagram{fig:m3properq}{三力与速度平行而功率非零}{\draw[axis](0,0)--(3.5,0)node[right]{$\vct v$};\draw[axis](0,.6)--(2.5,.6)node[right]{$\vct f$};}
\end{exercise}
\begin{solution}
图\ref{fig:m3properq}中，取瞬时速度沿$x$正向，平行三力的有向分量为$f$。四力为
\[
 F^\mu=\gamma_v(\beta_v f,f,0,0),\qquad
 F^2=\gamma_v^2f^2(\beta_v^2-1)=-f^2.
\]
故固有三力大小$f_0=\sqrt{-F^2}=|f|$。这里所谓四力模是类空四力的$\sqrt{-F^2}$，不是四个分量平方和的欧氏模。
\Diagram{fig:m3propera}{类空四力可有非零时间分量；零功率参考系中的时间分量为零}{\draw[axis](0,0)--(3.5,0)node[right]{$F^x$};\draw[axis](0,0)--(0,2.6)node[above]{$F^0$};\draw[->,thick](0,0)--(2.8,1.4);\draw[guide](0,0)--(2.4,2.4);\node[right]at(2.8,1.4){$F^0/F^x=v/c$};}
图\ref{fig:m3propera}中$F^0=\gamma_vfv/c$，只要$fv\ne0$便非零。变到粒子瞬时静止系后，
\[
 F_*^0=\gamma_v(F^0-\beta_vF^x)=0,\qquad F_*^x=f.
\]
但零时间分量并非瞬时静止系独有。以该静止系为起点，再沿$y$以任意$w<c$变换：原四力$(0,f,0,0)$的$0,y$分量均为零，变换后仍是$(0,f,0,0)$；粒子却以$-w\hat{\vct y}$运动。此系三力为$(f/\gamma_w,0,0)$，与粒子速度垂直，功率仍为零。

一般判据是$F^0=\gamma\vct f\cdot\vct v/c=0$，即该系中瞬时功率为零。四维正交$u\cdot F=0$则在所有参考系成立，它不要求每个参考系中的时间分量都为零。
\end{solution}

\begin{exercise}\label{ex:m3threshold}
静质量$m$的入射粒子撞击同样静质量的静止靶粒子，末态需产生四个质量均为$m$的粒子。忽略其他限制，求入射动能阈值，并与两束等能反向对撞比较。
\Diagram{fig:m3thresholdq}{固定靶与对撞的入射配置}{\draw[->](0,.5)--(2,.5);\fill(2.8,.5)circle(2pt);\node[above]at(1,.5){$E$};\node[right]at(3,.5){固定靶};\draw[->](0,-1)--(1.7,-1);\draw[->](3.6,-1)--(1.9,-1);\node[right]at(4,-1){对撞};}
\end{exercise}
\begin{solution}
在图\ref{fig:m3thresholdq}的固定靶实验室系，入射总能量记为$E$，动量记为$p$，有$E^2-p^2c^2=m^2c^4$。初态总能量$E+mc^2$、总动量$p$，因此
\begin{align*}
 s&=(E+mc^2)^2-p^2c^2\\
 &=E^2-p^2c^2+2mc^2E+m^2c^4\\
 &=2m^2c^4+2mc^2E.
\end{align*}
阈值时四个末态粒子在质心系相对静止，最小总能量为$4mc^2$，故$s=16m^2c^4$，解得$E=7mc^2$、$K=6mc^2$。
\Diagram{fig:m3thresholda}{阈值末态在质心系中无相对运动}{\foreach\x in{0,1,2,3}{\fill(\x,0)circle(3pt)node[below]{$m$};}\draw[curve](-.4,-.5)rectangle(3.4,.6);\node[above]at(1.5,.6){总静能$4mc^2$};}
图\ref{fig:m3thresholda}的末态静止是指质心系。实验室的初态总能量为$8mc^2$，总动量为$\sqrt{7^2-1}\,mc=4\sqrt3\,mc$，因而质心速度$V=Pc^2/E_T=\sqrt3c/2$、$\gamma_V=2$。末态整体动能为
\[
 K_{\rm bulk}=(\gamma_V-1)\,4mc^2=4mc^2.
\]
六份入射动能中，$2mc^2$用来增加末态静能，$4mc^2$用于整体运动。若改成两束等能反向对撞，总动量一开始为零，阈值只需每束总能量$2mc^2$，每束动能$mc^2$，两束共$2mc^2$。两种装置的末态最小静能相同，固定靶所需额外能量来自总动量约束。
\end{solution}

\begin{exercise}\label{ex:m3compton}
能量$E$的光子撞击静止电子，散射光子与入射方向成$\theta$角。用四动量守恒推导波长改变量。
\Diagram{fig:m3comptonq}{散射光子与反冲电子}{\draw[->,thick](-2,0)--(0,0)node[midway,above]{$E$};\draw[->,thick](0,0)--(2,1.7)node[right]{$E'$};\draw[->](0,0)--(2,-1.1)node[right]{电子};\draw[guide](0,0)--(2.5,0);\node at(.8,.35){$\theta$};}
\end{exercise}
\begin{solution}
图\ref{fig:m3comptonq}中，初始电子四动量为$P=(mc,\vct0)$，初末光子四动量记为$k,k'$，其能量分别为$E,E'$。末电子$P'=P+k-k'$，电子静质量不变，故
\[
 0=P'^2-P^2=2P\cdot k-2P\cdot k'-2k\cdot k'.
\]
光子各自平方为零，两个光子的空间夹角为$\theta$，因此
\[
 P\cdot k=mE,\quad P\cdot k'=mE',\quad
 k\cdot k'=\frac{EE'}{c^2}(1-\cos\theta).
\]
代入并除以$EE'$，得$1/E'-1/E=(1-\cos\theta)/(mc^2)$。
\Diagram{fig:m3comptona}{散射角越大，光子能量损失越大}{\draw[axis](0,0)--(4.1,0)node[right]{$\theta$};\draw[axis](0,0)--(0,2.8)node[above]{$\lambda'-\lambda$};\draw[curve,domain=0:180,samples=70]plot({\x/50},{1-cos(\x)});\node[below]at(3.6,0){$\pi$};\node[left]at(0,2){$2h/(mc)$};}
以$E=hc/\lambda$表示，即为图\ref{fig:m3comptona}的康普顿关系
\[
 \lambda'-\lambda=\frac h{mc}(1-\cos\theta),\qquad
 E'=\frac E{1+(E/mc^2)(1-\cos\theta)}.
\]
电子反冲不能在守恒中省掉。取散射光子在$x,y$平面且$y$分量为正，末电子分量为
\[
 p_{ex}=\frac{E-E'\cos\theta}c,\qquad
 p_{ey}=-\frac{E'\sin\theta}c,
 \quad K_e=E-E'.
\]
这些分量应与题图中的反冲方向一致。由它们平方得到$p_e^2c^2=E^2+E'^2-2EE'\cos\theta$；结合$K_e$又有$p_e^2c^2=K_e^2+2mc^2K_e$，两式相等恰恢复上面的散射能量式。正向散射无能量转移，反向散射波长改变量最大为$2h/(mc)$。
\end{solution}

\begin{exercise}\label{ex:m3circular}
相对论带电粒子在均匀磁场中作圆周运动，速度垂直磁场。忽略辐射。求半径、实验室角频率、固有加速度及瞬时静止系电场。
\Diagram{fig:m3circleq}{垂直磁场中的圆周运动，取正电荷}{\draw[curve](0,0)circle(1.5);\fill(1.5,0)circle(2pt);\draw[axis](1.5,0)--(1.5,1.2)node[right]{$v$};\draw[axis](1.5,0)--(.3,0)node[above]{$qvB$};\draw(0,0)circle(.15);\draw(-.1,-.1)--(.1,.1);\draw(-.1,.1)--(.1,-.1);\node[below]at(0,-.3){$B$入纸面};}
\end{exercise}
\begin{solution}
磁力始终垂直速度，$\dd E/\dd t=0$，故$\gamma_v$恒定。图\ref{fig:m3circleq}中动量转向率$p v/R=|q|vB$，因此$R=p/(|q|B)=\gamma_vmv/(|q|B)$，$\omega=v/R=|q|B/(\gamma_vm)$。
\Diagram{fig:m3circlea}{局部瞬时静止系中的电场指向圆心}{\fill(0,0)circle(2pt)node[below]{粒子};\draw[axis](0,0)--(-2,0)node[above]{$\vct E_0=\gamma_v\vct v\times\vct B$};\draw[guide](0,-1)--(0,1.2);\node[right]at(.2,.7){$\vct v$切向};}
实验室向心加速度$a=v^2/R$属垂直分量，所以固有加速度$\alpha=\gamma_v^2v^2/R=\gamma_v|q|vB/m$。瞬时静止系场变换给图\ref{fig:m3circlea}的$\vct E_0=\gamma_v\vct v\times\vct B$；静止粒子所受电力大小$|q|\gamma_vvB=m\alpha$，与加速度关系相符。该系电场可做瞬时加速，却不意味着实验室磁场做功。
\end{solution}

\begin{exercise}\label{ex:m3sheet}
无限带电平面静止系面密度为$\sigma_0$，平面为$xy$面。观察系沿$x$以$V$运动。求平面两侧的电磁场，并用变换后的面电流核对磁场方向。
\Diagram{fig:m3sheetq}{静止带电平面两侧的电场}{\draw[curve](-2,0)--(2,0);\foreach\x in{-1.5,-.5,.5,1.5}{\node at(\x,0){$+$};\draw[axis](\x,.2)--(\x,1.2);\draw[axis](\x,-.2)--(\x,-1.2);}\node[right]at(2,0){$xy$面};\node[right]at(1.5,1.2){$z>0$};}
\end{exercise}
\begin{solution}
图\ref{fig:m3sheetq}中$E_z=\operatorname{sgn}(z)\sigma_0/(2\varepsilon_0)$、$B=0$。沿$x$变换得$E_z'=\gamma E_z$、$B_y'=\gamma VE_z/c^2$。面密度$\sigma'=\gamma\sigma_0$，平面在新系沿$-x$运动，面电流$K_x'=-\gamma\sigma_0V$。
\Diagram{fig:m3sheeta}{新系面电流沿$-x$，上侧磁场沿$+y$}{\draw[axis](-2,0)--(2,0)node[right]{$x'$};\draw[axis](0,-1.7)--(0,1.8)node[above]{$z'$};\draw[very thick,->](1,0)--(-1,0)node[below]{$\vct K'$};\draw(.8,1)circle(.15);\draw(.7,.9)--(.9,1.1);\draw(.7,1.1)--(.9,.9);\node[right]at(1.1,1){$B_y'$入纸面};}
按图\ref{fig:m3sheeta}的右手定则，无限面电流上侧$B_y'=-\mu_0K_x'/2=\mu_0\gamma\sigma_0V/2$，与$\gamma VE_z/c^2$完全一致，因为$\varepsilon_0\mu_0c^2=1$。上下侧场都反号。这里$x'$向右、$z'$向上时，$+y'$朝纸内，基底方向需在图中明确。
\end{solution}

\begin{exercise}\label{ex:m3gauge}
真空平面波的四势为$A^\mu=\epsilon^\mu\ee^{-ik_\nu x^\nu}$，满足洛伦兹规范。求极化四矢量的约束，并说明沿波矢的四势分量为何不改变场。
\Diagram{fig:m3gaugeq}{横向极化与传播方向}{\draw[axis](0,0)--(4,0)node[right]{$\vct k$};\draw[axis](1,0)--(1,1.8)node[above]{$\vct\epsilon_\perp$};\draw[guide](0,1)--(4,1);\node[below]at(2,-.4){等相位波面垂直$\vct k$};}
\end{exercise}
\begin{solution}
规范条件$\partial_\mu A^\mu=0$给$k_\mu\epsilon^\mu=0$。若再选无标势的横向规范$\epsilon^0=0$，则$\vct k\cdot\vct\epsilon=0$，得到图\ref{fig:m3gaugeq}的两个独立横向极化方向。
\Diagram{fig:m3gaugea}{规范变化沿四维波矢方向，横向物理场不变}{\draw[axis](0,0)--(4.3,0)node[right]{$k^\mu$方向};\draw[axis](0,0)--(0,2.2)node[above]{横向};\draw[->,thick](0,0)--(1.5,1.6)node[above]{$\epsilon^\mu$};\draw[->,thick](0,0)--(3.2,1.6)node[right]{$\epsilon^\mu+ak^\mu$};\draw[guide](1.5,1.6)--(3.2,1.6);}
场张量由反对称组合$k^\mu\epsilon^\nu-k^\nu\epsilon^\mu$给出。把图\ref{fig:m3gaugea}中$\epsilon^\mu$改为$\epsilon^\mu+ak^\mu$，新增项$a(k^\mu k^\nu-k^\nu k^\mu)=0$，所以场不变。并且$k^2=0$使$k\cdot\epsilon$约束仍成立。图为四维分量关系示意，不是普通空间中把电场沿传播方向倾斜。
\end{solution}

\begin{exercise}\label{ex:m3parallelfields}
某处电磁场平行且均非零。能否通过惯性变换同时消去其中一个场？若两场本来不平行，两个场不变量可以提供哪些必要限制？
\Diagram{fig:m3parallelq}{平行电磁场的点积不为零}{\draw[axis](0,0)--(0,1.65)node[above]{$\vct E$};\draw[axis](1.2,0)--(1.2,1.35)node[above]{$\vct B$};\node[below]at(.6,0){同向或反向均可};}
\end{exercise}
\begin{solution}
图\ref{fig:m3parallelq}中$\vct E\cdot\vct B=\pm EB\ne0$，它是洛伦兹不变量。若在某系消去任意一个场，点积就为零，矛盾。因此不存在这种惯性系。
\Diagram{fig:m3parallela}{消场必须先满足两个不变量的限制}{\draw[axis](0,0)--(5,0)node[right]{$E^2-c^2B^2$};\draw[curve](2.5,-.6)--(2.5,.8);\node[below]at(2.5,-.6){$0$};\node[above]at(1,.2){纯磁场可能};\node[above]at(4,.2){纯电场可能};\node[below]at(2.5,-1.1){还须$\vct E\cdot\vct B=0$};}
一般场若可变为纯电场，必要有$\vct E\cdot\vct B=0$且$E^2-c^2B^2>0$；若可变为纯磁场，则后一不变量小于0。图\ref{fig:m3parallela}的零边界需要另行处理：两个不变量都零的非零场不能被消去，因为若某系变成零场，逆变换会使原系也为零。对正交场，例\ref{ex:m3crossed}已给出满足这些条件的具体速度。
\end{solution}

